% !TEX program = xelatex
\documentclass[11pt]{article}

\usepackage{fontspec}
\usepackage{unicode-math}
\setmainfont{Latin Modern Roman}[SmallCapsFont={Latin Modern Roman Caps}]
\setmathfont{Latin Modern Math}
\setmonofont{DejaVu Sans Mono}[Scale=0.82]
\usepackage[margin=1in]{geometry}
\usepackage{enumitem}
\usepackage{booktabs}
\usepackage{longtable}
\usepackage{array}
\usepackage{url}
\usepackage[hidelinks]{hyperref}

\DeclareUrlCommand\lean{\urlstyle{tt}}
\makeatletter
\g@addto@macro\UrlBreaks{\do\_\do\.}
\makeatother

\setlist{itemsep=2pt,topsep=3pt}
\setlength{\parindent}{0pt}
\setlength{\parskip}{4pt}
\setlength{\emergencystretch}{3em}
\Urlmuskip=0mu plus 1mu

% An item: #1 title, #2 status line, #3 counted false statements, #4 counted defective proofs.
\newcounter{defect}
\newenvironment{defect}[4]{%
  \par\addvspace{10pt}\refstepcounter{defect}%
  \noindent\textbf{Item \thedefect.\enspace #1}\par\nopagebreak
  \noindent\textit{#2}\par\nopagebreak\smallskip}{\par\addvspace{6pt}}
\newcommand{\fld}[1]{\par\noindent\textsc{#1.}\enspace}

\title{Defects in the published proofs of the classification\\ of homogeneous plane continua,
  ordered by severity}
\author{Working document compiled from the defects ledger of the formalization}
\date{22 October 2026}

\begin{document}
\maketitle

\begin{abstract}
\noindent The classification of homogeneous plane continua (Hoehn and Oversteegen, 2016) and the
papers it rests on have been formalized in Lean~4 with Mathlib. Every statement of the
classification, in its plane form and in the form with Bing and Jones's third class, is verified.
This document lists what the formalization found wrong or missing in the published texts, most
serious first. Two counts are kept: 23 printed statements are false, and 18 printed proofs fail at
an identified step while the statement is true. The counts measure kind, not weight. Fifteen of
the 23 false statements fail only at an empty set, a single point, an endpoint or a chain of at
most three links; in ten of the 18 failed proofs the argument stands once one step is repaired or a
small case is set aside. The defects that matter most are in neither count: the step from Jones's
decomposition to the circle of pseudo-arcs, used without proof or citation in the classification
paper, whose first proof found is in a 2022 preprint of Domokos and Prajs; and Bing and Jones's
sketch of a plane circle of pseudo-arcs.
\end{abstract}

\tableofcontents

%======================================================================
\section{Reader's guide}\label{sec:guide}

\textbf{What is listed.} Each item is a defect of a published text, found while formalizing it.
Items are numbered consecutively through the document. Each item gives the source and page, the
printed text where it matters, what fails, how the formalization settles it, and a status line
in italics. The status line names the label from the list below and the section of the project's
defects ledger (``Ledger A.23'') where the entry is recorded in full.

\textbf{Labels.}
\begin{description}[leftmargin=1.5em,style=nextline]
\item[False as printed] The printed statement is refuted by a verified counterexample.
\item[Proof defective; statement true] The printed proof fails at an identified step, and the
  statement is proved by a repaired or different argument.
\item[Not established] An assertion of a proof sketch that the formalization neither proves in the
  generality stated nor refutes.
\item[Unstated step] The proof uses a fact the paper neither states nor proves. The fact is true
  and is proved in the formalization.
\item[Unargued] A sentence of a proof is asserted without argument. It is supplied, or shown
  unnecessary, in the formalization.
\item[Stated without proof] No proof is printed.
\item[Citation] A cited source does not contain, or does not prove, what it is cited for.
\item[Harmless] Wording that disagrees with itself or with a cited source, without consequence.
\end{description}
Only the first two labels are counted. A defective proof is counted only when its statement is
settled; Oversteegen--Tymchatyn's Theorem~3.2 (Item~\ref{it:ot32}) is therefore not counted. A
refuted printed sentence is not counted when the same paper states it correctly elsewhere
(Hoehn--Oversteegen's restatement of Theorem~1 on p.~211) or when it is a remark and not a
numbered statement (their remarks of pp.~193--194); both are listed in \S\ref{sec:wording} and
\S\ref{sec:degenerate}. A tacit hypothesis is counted only when a counterexample to the printed
form is verified.

\textbf{Verified.} A Lean theorem is \emph{verified} when it compiles and its axiom audit lists
only \lean{propext}, \lean{Classical.choice} and \lean{Quot.sound}. Every Lean name in this
document is listed in Appendix~\ref{app:lean} with its namespace and the axioms it depends on, as
recorded in the repository's \lean{axioms.txt} (4030 lines). A \emph{hand argument} is an argument
written out and read, but not formalized; it is marked as such wherever it occurs.

\textbf{Quotations} are transcribed from page images rendered at 200 dots per inch or more
(400 for subscripts, primes and bars), not from the text layer of the PDF files. Page numbers are
journal pages. Printed notation is kept: in the older papers $A+B$ is the union and $A\cdot B$ the
intersection.

\textbf{Order.} The sections run from the most serious defects to the most trivial:
\begin{center}
\begin{tabular}{@{}llcc@{}}
\toprule
Section & Kind & False & Proofs \\
\midrule
\ref{sec:path} & Gaps on the path of the classification & 0 & 5\\
\ref{sec:intended} & Statements false in configurations their authors meant to cover & 3 & 1\\
\ref{sec:other} & Other proofs that do not go through & 0 & 2\\
\ref{sec:cite} & Citations that do not carry what they are cited for & 0 & 0\\
\ref{sec:local} & Proof steps that fail locally & 0 & 10\\
\ref{sec:wording} & Statements false because of their wording & 5 & 0\\
\ref{sec:degenerate} & Statements false at degenerate cases & 15 & 0\\
\ref{sec:unargued} & Unargued steps and statements without proof & 0 & 0\\
\ref{sec:misprints} & Misprints and harmless mismatches & 0 & 0\\
\midrule
 & Total & 23 & 18\\
\bottomrule
\end{tabular}
\end{center}
Section~\ref{sec:path} concerns the published proofs of the results the classification is
assembled from: Hoehn and Oversteegen's Theorems~A, B and C and their Theorem~2, the arc step that
the proof of Theorem~B needs, the third class, and Bing's construction of the pseudo-arc.
Section~\ref{sec:other} concerns supporting results of plane topology used inside those proofs.
Moore 1919 Theorem~D is of the second kind; it is placed in \S\ref{sec:intended} beside Theorem~C,
whose false case its proof uses. Within a section, items are ordered by their consequence for the
classification. Appendix~\ref{app:readings} lists readings of printed statements that the
formalization refuted and that are \emph{not} defects of the sources, among them Bing 1951
Theorem~8 (Item~\ref{it:b51-8}).

%======================================================================
\section{The classification and what it rests on}\label{sec:classification}

Hoehn and Oversteegen \cite[pp.~179--180]{HO16} assemble the classification from four results:
\begin{quote}
\textsc{Theorem A.} ([16]) \emph{If $M$ is a decomposable homogeneous continuum in the plane,
then $M$ is a circle of mutually homeomorphic indecomposable homogeneous continua.}

\textsc{Theorem B.} ([12]) \emph{Every indecomposable homogeneous plane continuum is hereditarily
indecomposable.}

\textsc{Theorem C.} ([45]) \emph{Every homogeneous indecomposable plane continuum has span zero.}

\textsc{Theorem 1.} \emph{A non-degenerate continuum $X$ is homeomorphic to the pseudo-arc if and
only if $X$ is hereditarily indecomposable and has span zero.}
\end{quote}
Here [16] is Jones 1955 \cite{J55}, [12] is Hagopian 1976 \cite{H76}, and [45] is Oversteegen and
Tymchatyn 1982 \cite{OT82}. Theorem~1 is the paper's own. Then \cite[p.~180]{HO16}:
\begin{quote}
It follows immediately from Theorems B and C above and Theorem 1 that every indecomposable
non-degenerate homogeneous plane continuum is a pseudo-arc. Combining this with Theorem A above,
we obtain the following classification of homogeneous plane continua.

\textsc{Theorem 2.} \emph{Up to homeomorphism, the only non-degenerate homogeneous continua in the
plane are (1) the circle, (2) the pseudo-arc, and (3) the circle of pseudo-arcs.}
\end{quote}

\textbf{The formal statement.} \lean{ho16_classification_plane₄}: there are plane sets $C_1$,
$C_2$, $C_3$, homeomorphic respectively to the circle, a pseudo-arc and a circle of pseudo-arcs in
Bing and Jones's sense (a circle-like continuum with a continuous decomposition into pseudo-arcs
whose decomposition space is a simple closed curve), each homogeneous, no two homeomorphic, such
that a compact connected plane set with more than one point is homogeneous if and only if it is
homeomorphic to one of them. \lean{ho16_classification₄} states the same for a continuum given as
a subspace of the plane.

\textbf{Where each source enters.} The formalization reads each cited result in its source and
proves it, or proves what the classification needs from it.
\begin{itemize}
\item Theorem A: Jones 1955 Theorems~1 and 2 \cite{J55}, with Jones 1948, 1949, 1951
  \cite{J48,J49,J51}, Whyburn 1939 \S14 and Whyburn's \emph{Analytic Topology} \cite{W39,W42},
  Moore 1922 \cite{M22}, Moore 1925 \cite{M25} for the plane quotient, and the Effros theorem in
  van Mill's form \cite{vM04}. Items~\ref{it:j55}, \ref{it:w39}, \S\ref{sec:local},
  \S\ref{sec:degenerate}.
\item Theorem B: Hagopian 1976 \cite{H76}, with Hagopian 1975 \cite{H75}. The formalization
  proves it by Hagopian 1984 \cite{H84} and the arc step of Bing 1960 \cite{B60}, which it proves by
  a different argument from Jones 1975's \cite{J75}. Items~\ref{it:h76}, \ref{it:j75}.
\item Theorem C: Oversteegen and Tymchatyn 1982 Theorem~3.6, with Cook 1967 \cite{C67}.
  Item~\ref{it:ot32}.
\item Theorem 1: Hoehn and Oversteegen 2016 \cite{HO16}, with Bing's characterization of the
  pseudo-arc \cite{B48,B51a,B51b} and its supports \cite{L61,F66,F69,I68}. Theorem~1 is formalized
  in the direction and generality the classification uses: a hereditarily indecomposable plane
  continuum of span zero is snake-like (\lean{ho16_theorem1_plane}, whose hypothesis is
  Hoehn--Oversteegen's Theorem~20, verified as \lean{ho16_theorem20}); with Theorems~B and C, a
  nondegenerate plane continuum is homogeneous and indecomposable if and only if it is a pseudo-arc
  (\lean{ho16_indec_homogeneous_iff}). Items~\ref{it:b48-3}, \ref{it:b48-8}, and
  Hoehn--Oversteegen's own entries in \S\S\ref{sec:cite}--\ref{sec:misprints}.
\item The third class: Bing and Jones 1959 \cite{BJ59} and Domokos and Prajs 2022 \cite{DP22}.
  Items~\ref{it:unique}, \ref{it:bj7}, \ref{it:bj5}.
\item Plane topology used along the way: Moore 1916, 1919, 1920, 1925 \cite{M16,M19,M20,M25} and
  Mullikin 1922 \cite{Mu22}. Items~\ref{it:m19}, \ref{it:m25-21}, \S\S\ref{sec:local}--\ref{sec:degenerate}.
\end{itemize}
\textbf{What the final proof uses.} The repository's dependency census lists, declaration by
declaration, what the two final theorems use: 4689 of the library's 11697 declarations, in 446 of
its 863 modules. Their closure, 12480 constants of the library, is accepted by the kernel when it
is replayed into an environment built from Mathlib alone. Appendix~\ref{app:lean} marks each
declaration cited here that the final theorems use. The formal route differs from the published
one in places. It does not use Bing 1948 Theorems 3, 4, 5, 7, 8, 9 or 13--15, Bing 1951 (Duke) Theorems 8, 9,
11 and 12, Moore 1922, Jones 1948 Theorems 6 and 10, or Hagopian 1976's own proof of Theorem~1,
which it replaces by Hagopian 1984 (Items~\ref{it:h76}, \ref{it:b48-3}, \ref{it:b48-8},
\ref{it:b48-4}). The items are ordered by their place in the published argument.

%======================================================================
\section{Gaps on the path of the classification}\label{sec:path}

Each item in this section concerns the published proof of one of the results the classification
is assembled from (\S\ref{sec:guide}). In each, the printed argument is missing or cannot be
completed as written, and the result is proved from another source or by a new argument.

%----------------------------------------------------------------------
\begin{defect}{The step from Jones's decomposition to the circle of pseudo-arcs (Hoehn--Oversteegen
2016, Theorem~2)}{Unstated step in the proof of a main theorem; no proof found before
Domokos and Prajs's 2022 preprint. Not counted: no printed statement is false. Ledger A.25, A.26,
A.27.}{0}{0}\label{it:unique}

\fld{Printed} Theorem~2 is obtained by ``Combining this with Theorem A above'' (\S\ref{sec:classification}).
The circle of pseudo-arcs is defined on pp.~182--183 as a continuum that ``is a circle-like
continuum which admits an open map to the circle whose point preimages are pseudo-arcs'', and
p.~183 adds: ``Bing and Jones proved that the circle of pseudo-arcs is homogeneous and that it is
unique, in the sense that it is the only continuum (up to homeomorphism) with the above
properties.''

\fld{What is missing} For a decomposable homogeneous plane continuum $M$ that is not a simple
closed curve, Theorem~A with Theorems~B, C and 1 gives an open map of $M$ onto the circle whose
point preimages are pseudo-arcs
(p.~179: ``The conclusion of this theorem implies that there is an indecomposable homogeneous
continuum $X$ (possibly a single point) and an open map from $M$ to the circle all of whose point
preimages are homeomorphic to $X$''). Bing and Jones's uniqueness applies to \emph{circle-like}
continua: their definition (p.~173) reads ``It is a circle-like continuum $M$ such that there is a
continuous decomposition of $M$ into pseudo-arcs such that the decomposition space is a simple
closed curve.'' That $M$ is circle-like is used and is neither stated nor cited.

\fld{Why it needs an argument} An open map onto the circle with pseudo-arc fibres does not give
circle-likeness. The product of the pseudo-arc with the circle has such a map and is
two-dimensional, so not circle-like; Hoehn and Oversteegen mention it (p.~183: ``homogeneous but
not embeddable in the plane''). Planarity and dimension one do not suffice either. Prajs
\cite[Theorem~8, p.~1753]{P00}: ``There exist a plane s-arc of pseudo-arcs and a plane s-circle
of pseudo-arcs'', curves with continuous decompositions into pseudo-arcs over an arc and over a
simple closed curve; on the same page: ``the examples constructed here are triodic, and thus
topologically different. Evidently, their elements of the decompositions are not terminal.''
Circle-like continua are atriodic, so Prajs's plane s-circle is not circle-like. A proof must
therefore use homogeneity. What homogeneity supplies is that Jones's elements are terminal: Jones
1955 Theorem~1(c), p.~736: ``if $g$ is an element of $G$ and $K$ is a subcontinuum of $M$ which
contains both a point of $g$ and a point of $M-g$, then $g$ is a subset of $K$'' (verified,
\lean{jones55_c}).

\fld{The literature before 2022} Lewis \cite[p.~80]{Lw83} asks, as ``Question [28]'', whether two
(homogeneous) one-dimensional continua with continuous decompositions into pseudo-arcs over
homeomorphic quotients must be homeomorphic, and adds: ``This question has an affirmative answer if
all of the decomposition elements in each of $G_1$ and $G_2$ are terminal continua.'' His [28] is
the paper published as Lewis 1985, then submitted; p.~80 gives no proof.
Lewis 1985 \cite[p.~93]{Lw85}, the paper that constructs such continua with terminal elements:
``It will follow from our construction that all decomposition elements are terminal continua. We
still do not know in general whether $\hat M$ is unique.'' Lewis 1996 \cite[p.~173]{Lw96} combines
an observation of his own with Lewis 1985 (cited there as Theorem~2 [24], an existence and lifting
theorem) and concludes: ``For each continuum in one class there is a unique corresponding continuum
in the other class, obtained either by shrinking out maximal terminal pseudo-arcs or by enlarging
each point to a terminal pseudo-arc.'' Uniqueness in the enlarging direction is the question Lewis
1985 leaves open. It is not proved in any of the three papers.

\fld{The first proof found} Domokos and Prajs \cite[Theorem~6.1]{DP22}: ``A continuum $X$ is a
circle of pseudo-arcs if and only if $X$ admits an open atomic map $f\colon X\to S$ onto a simple
closed curve $S$ such that all point-inverses, $f^{-1}(y)$, are pseudo-arcs.'' Atomic means that
every point-inverse is terminal; ``circle of pseudo-arcs'' means homeomorphic to Bing and Jones's
example. The authors: ``We include the proof because we did not find in the literature a proper
reference to this theorem.'' The paper is an arXiv preprint; no journal version and no earlier
proof were found in a search of the papers on homogeneous continua available to the project. Their
proof shows that $X$ is atriodic (``one can easily show''; unargued, and not needed by the formal
route below), that the preimages of the two halves of $S$ are atriodic folders of chainable
continua, chainable by \cite[Proposition~6]{HMP20}, and glues them to Bing and Jones's example by
Bing--Jones Theorem~10.

\fld{Formalization} Verified, with no named hypothesis: Theorem~6.1 in both directions
(\lean{dp22_theorem61'}, \lean{dp22_theorem61_iff'}), the circle-likeness step
(\lean{hoCircleLike'}), and the classification (\lean{ho16_classification₄},
\lean{ho16_classification_plane₄}). The formal route makes the halves chainable by Bing's argument
for Theorem~11 of \cite{B51a} (\lean{isSnakeLike_of_atomic'}), which needs neither folders nor
shape theory, and glues them by Bing--Jones Theorem~10 (\lean{bjThm10}). Jones's decomposition of
$M$ supplies the atomic map through \lean{jones55_c}. In the formalization ``circle of
pseudo-arcs'' has Bing and Jones's meaning (the definition \texttt{IsCircleOfPA}: circle-like,
with an open map onto the circle whose fibres are pseudo-arcs). Domokos and Prajs's meaning,
homeomorphic to Bing and Jones's example, is not formalized, since that example is not constructed;
by the two theorems above, any two circles of pseudo-arcs in Bing and Jones's sense are
homeomorphic.
\end{defect}

%----------------------------------------------------------------------
\begin{defect}{A plane circle of pseudo-arcs (Bing--Jones 1959, \S7, pp.~189--191)}{Not
established: a sketch whose key step is not argued. Not counted: no instance shows the step
failing. The claim is verified by a different construction. Ledger A.25 (sixth reading).}{0}{0}\label{it:bj7}

\fld{Printed} \S7 begins (p.~189): ``It follows from Theorem 10 that each circle of pseudo-arcs
is homogeneous and that any two of them are homeomorphic. In this section we show that the plane
contains a circle of pseudo-arcs.'' On p.~190: ``Now there exists a sequence $D_1, D_2, D_3,\cdots$ of circular chains
of simple domains in the plane such that'' conditions (1)--(9) hold. On p.~191: ``With the
addition of (9) one may see with the help of [3] that each element of $G$ is a pseudo-arc. It
follows from Theorem 10 that $M$ is homogeneous.''

\fld{What is missing}
\begin{itemize}
\item The chains $D_i$ are asserted to exist and are not constructed.
\item That the elements of $G$ are disjoint, fill $M$, and form a continuous decomposition whose
  quotient is a circle is not argued, with or without (9). Unargued also: ``Also, in this case $M$
  would be a continuous circle-like circle of continua'' (p.~191, of the chains without (9)).
\item Not established: that each element of $G$ is a pseudo-arc. Condition (9) makes each new
  column $T_{ij}$ crooked in a pair $T_{hk}+T_{hl}$ of old columns. It does not control an
  element of $G$ covered at some level by two adjacent columns, and two crooked chains joined end
  to end are not crooked in general. No chains satisfying (1)--(9) with an element of $G$ that is
  not a pseudo-arc are known, so the step is not shown to fail.
\end{itemize}

\fld{Consequence} The third class of the classification is a \emph{plane} continuum only through
\S7. Hoehn and Oversteegen (p.~179): ``Bing and Jones [5] constructed in 1954 such a continuum in
the plane for which $X$ is the pseudo-arc, and also proved that it is homogeneous.'' Lewis
\cite[p.~91]{Lw85} also takes \S7 at face value. His own construction of a circle of pseudo-arcs
\cite[p.~93]{Lw85} identifies ``the end pseudo-arcs of arcs of pseudo-arcs in the same manner that
the endpoints of arcs are identified''; it takes the arc of pseudo-arcs from Bing and Jones and does
not discuss planarity or circle-likeness.

\fld{Formalization} Verified: \lean{bjPlaneExample}, through \lean{exists_circDescr_plane} and
\lean{CircDescr.isCircleOfPA}. The construction is not Bing and Jones's. It uses nested circular
chains of plane domains whose links are neighbourhoods of arcs of a base circle; the pattern of
each chain in the previous one is the lift of a degree-one map built from \emph{sweeps}, along each
of which the window of links changes at one end at a time (\lean{sweep_lemma},
\lean{sweep_lemma3}); and a planar step (\lean{circ_step}) obtained by descending an equivariant
shear of a strip through the exponential map. A recognition theorem shows that the intersection is
a circle-like continuum with a continuous decomposition into pseudo-arcs over the circle
(\lean{CircDescr.isCircleLike_M}). Nothing of \S7 is used except its claim.
\end{defect}

%----------------------------------------------------------------------
\begin{defect}{Hagopian 1976, Theorem~1 (Theorem B), pp.~345--349}{Proof defective; statement true.
Counted (one defective proof). The printed route is not completed; the theorem is verified by
Hagopian's 1984 route. Ledger A.23, A.24.}{0}{1}\label{it:h76}

\fld{Statement} Every indecomposable homogeneous plane continuum is hereditarily indecomposable.
The proof argues by contradiction from a type~A$'$ subcontinuum $E$, through an induction over the
blocks $x_1<x_2<\cdots<x_{2n}$ of a continuum $F$ in a strip $\pi^{-1}[D]$, ``Following Bing [2,
p.~225]'' (p.~348), that is, Bing 1960 \cite[pp.~225--227]{B60}.

\fld{The failing step} p.~348: ``The proof now breaks down into two cases''. The case split
presupposes that $d^{-1}(x_{2i-1})$ and $d^{-1}(x_{2i-2})$ each meet a side segment
$\pi^{-1}[\langle p_1,r\rangle]$ or $\pi^{-1}[\langle p_5,s\rangle]$. At the last step of the
induction ($i=2$, $J=[x_1,x_2]=[0,e]$) this is false: $d^{-1}(x_2)=d^{-1}(e)=k^{-1}(1)$ lies over
the far end of $\pi[E]$ and meets neither side segment. So neither printed case applies to the
step that produces the final contradiction with ``$[x_1,x_2]$ does not have Property 4''.
Certificate: \lean{blockLayersMeetSides_false}. In Bing's model \cite[p.~226]{B60} the
corresponding points are end points of arcs irreducible between the two side lines, so they lie on
the sides by construction; Hagopian's block ends are values of $d$, and nothing printed places their
layers on the sides. For $i\ge3$ the presupposition is true (\lean{blockLayersMeetSides3}), and the
last step can be read separately with the bottom edge in place of the strand
(\lean{last_step_ending}).

\fld{A lemma from the 1984 paper} ``Note that $x_1=0$, $x_2=e$, and $x_{2n}=1$''
(pp.~347--348), in its last clause, and Property~4 for the top
block need a saturation property the text does not state: every component of $F\cap\pi^{-1}[D]$
through a point over the inner middle of $D$ crosses $D$. The formalization proves it
(\lean{innerComponentsCross_hagopian}), and at the end layer $d^{-1}(1)$ of $F$ only with Hagopian
1984 Lemma~7 (\lean{hagopian84_lemma7}): when two type $\lambda$ continua overlap without either
containing the other, every layer of one that meets the other is a layer of the other. The 1976
paper has nothing of this kind at end layers; its Lemma~3 covers interior layers.

\fld{Printed sentences neither proved nor refuted} By a verified theorem
(\lean{hagopian76_theorem1_of_cases4}) the 1976 route reduces to four named claims, which are
open. They collect the following sentences of pp.~348--349.
\begin{itemize}
\item Case 1: ``By considering the position of $d^{-1}[[x_{2i-3},x_{2i}]]$, we see that $y_{i-1}$
  is between $d^{-1}[I]$ and $d^{-1}[J]$''. The position argument needs the unicoherence of the
  plane, which is neither stated nor cited; Hagopian's strands, components of
  $F\cap\pi^{-1}[D]$, need not end on the side lines as Bing's arcs do, and $L$ may re-enter
  $\pi^{-1}[D]$ on the far side of a strand.
\item Case 1: the segment $Q$ of the vertical over $p_4$ closes the band between two blocks only if
  the band crosses that vertical once.
\item p.~348: ``the continuum $N$ does not come within $\varepsilon_2$ of $\pi^{-1}(p_4)$''. Points
  of $N$ outside $\pi^{-1}[D]$ are not addressed. Bing's corresponding sentence \cite[p.~227]{B60}
  is unargued as well.
\item Case 2 (p.~349): ``If $y_{i-1}$ is above $d^{-1}[J]$, it is between $d^{-1}[I]$ and
  $d^{-1}[J]$, and $f(y_{i-1})$ is above $d^{-1}[I]$'' and ``Hence $y_{i-1}$ is below $d^{-1}[J]$
  and $f(y_{i-1})$ is above''. No reason is given. Bing's corresponding sentence follows from the
  parity of crossings of a simple closed curve (a hand argument on polygons); for Hagopian's
  continua there is no parity argument.
\item p.~348: ``(Other cases are handled with similar arguments to that given in this case.)'' Of
  the four orientations, two are treated by the formalization; in the other two the frame is not
  symmetric (the bottom edge lies over $E$, the top over $I_2$).
\item p.~348, E2.4: ``The point $y_{i-1}$ is below $C_i$ and $\rho(y_{i-1},d^{-1}[I])\ge
  \varepsilon_3$; for otherwise, a subcontinuum of $L$ would violate Property 2'' is verified for
  the strand next to $C_i$ only, conditionally on three sub-claims; the printed form, for all of
  $d^{-1}[I]$, is not decided.
\item p.~349, T17: ``Hence $f(y_{i-1})$ is either above $d^{-1}[I]$ or below $d^{-1}[J]$''
  presupposes that every point off the crossing continuum is above or below it. Components of
  $F\cap\pi^{-1}[D]$ that meet the frontier only in a side segment are not addressed.
\end{itemize}

\fld{Formalization} Verified: \lean{hagopian76_theorem1}, by Hagopian 1984 Corollary~3, which
carries out F.~B.~Jones's plan: Hagopian 1984 Theorem~1 with $N=M$ (\lean{hagopian84_decomp})
gives a monotone continuous decomposition with a homogeneous plane quotient containing an arc, and
the arc step (Item~\ref{it:j75}) gives the contradiction. About twenty further unargued steps of
pp.~340--349 are supplied (\S\ref{sec:unargued}).
\end{defect}

%----------------------------------------------------------------------
\begin{defect}{The arc step: Jones 1975, property (f) and the Theorem, pp.~59--61}{Two proofs
defective; statements true, that of (f) only vacuously in the plane. Counted (two defective
proofs). Verified by a different argument.
Ledger A.28.}{0}{2}\label{it:j75}

\fld{Role} Bing 1960 \cite{B60} proves that a homogeneous bounded plane continuum containing an arc
is a simple closed curve. Hagopian's 1984 route to Theorem~B needs it for indecomposable continua.
Jones 1975 \cite{J75} gives a short proof: ``\textsc{Theorem}. \emph{If $M$ is a homogeneous plane
continuum containing an arc, then $M$ is a simple closed curve.}'' (p.~59).

\fld{Property (f), p.~59} ``If $p$ is a point of $M$ and $D$ is an open set in $E$ with
$p\notin\bar D$, then there exists a positive number $\epsilon$ such that no arc $abc$ in $M$ has
the property that both $a$ and $c$ belong to $D$ and the $\epsilon$-neighbourhood of the subarc
$ab$ contains both $p$ and the subarc $bc$.'' The proof ends: ``Since $b'c'$ is in the
$\epsilon$-neighbourhood of $a'b'$, the image of $b'c'$ must be in the $3\epsilon$-neighbourhood of
the image of $a'b'$. This is impossible since the $3\epsilon$-neighbourhood of $q_1p_1$ contains no
point of $p_2q_2$ (and conversely).'' This does not follow: the image of $a'b'$ may run along the
arc $q_1pq_2$ past $q_1$ and come back near $p_2q_2$. The proof uses only the Effros theorem and
property (c), hence only hypotheses that the dyadic solenoid satisfies (homogeneous,
indecomposable, hereditarily unicoherent, atriodic, with (c)); in the solenoid (f) fails, since a
long arc of a dense line is $\epsilon$-dense (a hand argument, not formalized). Two earlier steps
are unargued as well: ``takes $b$ to $p$'' and ``$b'$ is close enough to $b$'' need $b$ near $p$,
and ``leaves the images $q_1$ and $q_2$ of $a$ and $c$ respectively in $D$'' is unjustified for a
homeomorphism that is not small. For plane continua (f) holds only vacuously, once the Theorem
is known (no Lean declaration states (f)).

\fld{The Theorem, pp.~59--61} The step ``This makes $h$ take $[t_0,t_d]$ onto (approximately)
$[t_d,t_0]$. So there is a fixed point $t_f$'' is false for small $\epsilon$: a small
homeomorphism preserves the orientation of the ray $T$ (\lean{JonesSetting.orientation}, verified).
Under the hypotheses of the setting the step contradicts this orientation lemma, so any argument for
it would itself finish the proof. The Theorem also uses (f) twice, and four further steps are
unargued (\S\ref{sec:unargued}).

\fld{Formalization} Verified: \lean{jones75_noArc} (no indecomposable homogeneous plane continuum
contains an arc) and \lean{bing60_plane_uncond} (Jones's Theorem). The proof is a
Poincaré--Bendixson argument in a collar box around an arc of a line, not Jones's: by Effros and
the orientation lemma, points of the line near the base arc lie on left-to-right passes of the box;
the region between the base pass and a later thin pass is bounded by a Jordan curve made of a
vertical segment and an arc of the line; once the line leaves that region it never re-enters it
(\lean{trap_jordan}), which contradicts recurrence (\lean{jones75_contra}). Jones's (d) is
verified in both directions (\lean{JonesSetting.jones75_d}), and (c) through
\lean{JonesSetting.union_arc} and \lean{JonesSetting.nested}.
\end{defect}

%----------------------------------------------------------------------
\begin{defect}{Jones 1955, Theorem~1 (Theorem A), the case argument of pp.~737--738}{Proof
defective; statement true. Counted (one defective proof). Ledger A.11.}{0}{1}\label{it:j55}

\fld{Statement} p.~736: ``\textsc{Theorem 1.} \emph{Suppose that $M$ is a decomposable,
homogeneous, compact, metric continuum. Then there exists a nondegenerate continuous collection
$G$ of mutually exclusive continua filling up $M$ such that (a) with respect to its elements (as
points) $G$ is a homogeneous, aposyndetic, compact, metric continuum, (b) if $x$ is a point of $M$,
then $L_x$ is an element of $G$, and (c) if $g$ is an element of $G$ and $K$ is a subcontinuum of
$M$ which contains both a point of $g$ and a point of $M-g$, then $g$ is a subset of $K$.}'' Here
$L_x$ is $x$ together with the points $y$ at which $M$ is not aposyndetic with respect to $x$, and
$U_x=M-L_x$.

\fld{What fails} In Case~I, ``Since $M_0$ is a subset of $M_1$'' fails as written: when $o(q)\in
V_2$, the irreducible continuum $K_o(q)$ can be $\{o\}$, which is not in $M_1$; the text does not
keep $o(q)$ out of $V_2$. And ``Since $q$ does not cut $o$ from $U_q$'' draws on $q\notin N_o$,
which is stated with $U_o$; it therefore needs $U_q=U_o$, which is (b), not yet available at that
point. The argument serves $N_o=M-U_o$.

\fld{Formalization} Verified: (a), (b), (c), upper semicontinuity, continuity, and that the
elements are homogeneous and mutually homeomorphic (\lean{jones55_a}, \lean{jones55_b},
\lean{jones55_c}, \lean{jones55_usc}, \lean{jones55_continuous}, \lean{jones55_elements}). Jones's
claims (1)--(10) and his properties (a)--(f) of pp.~738--739 are kept. The case argument and the
upper-semicontinuity argument are replaced by a spreading lemma from the Effros theorem: in a
homogeneous compact metric space, a continuum with interior that contains $o$ and misses $x$ gives
aposyndesis at $o$ with respect to $x$ (\lean{aposAt_of_isContinuum_interior}).
\end{defect}

%----------------------------------------------------------------------
\begin{defect}{Bing 1948, Theorem~3, the special case, p.~732}{Proof defective; statement true.
Counted (one defective proof). Ledger A.1.1.}{0}{1}\label{it:b48-3}

\fld{Statement} ``\textsc{Theorem 3.} \emph{If the chain $D$ is crooked in the chain
$E=[e_1,e_2,\cdots,e_m]$ and $d_i$ is a particular link of $D$, then there is a chain $F$ such
that $F$ is a consolidation of $D$, $d_i$ is contained in only the first link of $F$, each link of
$F$ is a subset of the sum of two adjacent links of $E$, and any link of $F$ containing an end link
of $D$ which intersects $e_1$ or $e_m$ is a subset of $e_1+e_2$ or $e_{m-1}+e_m$.}''

\fld{What fails} In the special case (no interior link of $D$ meets $e_1$ or $e_r$), with $d_i$ a
link of $D(1,h)$, Bing takes $H$ from the induction hypothesis and sets $F$ to be $H(1,u-1)$,
followed by the level groups of the remaining links $D'$, followed by $D(k+1,n)$. Certificate
(\lean{Recon121.seam_every_H_fails}, with \lean{Recon121.seamD_special}): eleven links of a
realized chain over five coarse links satisfy the special-case hypotheses, with $(h,k)=(4,6)$
forced; a chain $H$ satisfying the induction hypothesis exists; yet for every $i\le h$ and every
such $H$, the chain Bing writes down is not a chain with $d_i$ in its first link only.

\fld{Formalization} Verified: the theorem as printed in every clause
(\lean{theorem3_literal}, \lean{Recon121.theorem3LiteralOr_holds}, \lean{theorem3_bing}). The repair
(\lean{Recon122.bing_main}) makes three changes, none in the text: the level groups are anchored
at the bottom; an explicit grouping is made when there are five coarse links
(\lean{Recon122.special_five}); and the branch $d_i\in D(h,k)$ is handled by the induction
hypothesis applied to $D(h,k)$ inside $E(2,r-1)$.
\end{defect}

%----------------------------------------------------------------------
\begin{defect}{Oversteegen--Tymchatyn 1982, Theorem~3.2, pp.~137--138, in the proof of Theorem~C}{Proof defective; statement
not decided. Not counted. Ledger A.29.}{0}{0}\label{it:ot32}

\fld{Printed} The proof begins: ``We may suppose that for each $\lambda\in\Lambda$
$h_\lambda(Y)$ meets $\{0\}\times R$ in a single point and $h_\lambda(Y)$ intersects $\{1\}\times R$
in at most one point. Hence, \ldots\ we may suppose that for each $\lambda\in\Lambda$ neither
$h_\lambda(X)$ nor $h_\lambda(Y)$ separates the plane. Let $\delta\in\Lambda$. We may suppose that
$\pi_1(h_\delta(y))>\pi_1(h_\delta(z))+1/2$ for each $z\in X$. We may also suppose that
$d(h_\lambda(z),h_\gamma(z))<1/8$ for each $z\in Y$ and $\lambda,\gamma\in\Lambda$. The proof now
proceeds as in Theorem 3.1''.

\fld{What fails} The two suppositions about $\delta$ are not available together: the cluster of
maps within $1/8$ need not contain $\delta$, and one rescaling does not fix the other gaps. The
single-point supposition is unargued and can be avoided but not justified (disjoint crossing arcs
can each meet a side in two points). The proof of Theorem~3.1 picks three images one of which
separates the other two in the strip, which is unargued.

\fld{Formalization} What Theorem~C needs is a strip form in which the images lie in disjoint
continua crossing the strip, with their traces on the side ordered. That form is verified
(\lean{ot82_theorem32}). From it, Corollary~3.5 and Theorem~3.6 in the indecomposable case follow
(\lean{ot82_cor35_of_thm32}, \lean{ot82_theorem36_of_cor35}, two implications), and Theorem~C is
verified with no hypothesis (\lean{theoremC}). Whether the printed Theorem~3.2 holds in general is
not decided.
\end{defect}

%======================================================================
\section{Statements false in configurations their authors meant to cover}\label{sec:intended}

%----------------------------------------------------------------------
\begin{defect}{Bing--Jones 1959, Theorem~5, pp.~175--176}{False as printed. Counted (one false
statement). Its one use is existential and survives. Ledger A.25 (second reading).}{1}{0}\label{it:bj5}

\fld{Printed} ``Then if $U$ is an open set such that $U\cdot G\neq0$, $U\cdot(g_0+g_1)=0$, and
$\bar U$ lies in a link of $E$, then $G-U$ may be expressed as the sum of two mutually separated
sets $A$, $B$ such that $g_0\subset A$, $g_1\subset B$, and for each $x$ ($0\le x\le1$), either
$A\cdot g_x$ intersects each link of $E$ or $B\cdot g_x$ intersects each link of $E$.''

\fld{What fails} In every instance with $n\ge3$ links. Merge the last two links (replace
$e_{n-1}$ by $e_{n-1}\cup e_n$ and keep $e_n$), so that $\overline{e_n}\subseteq e_{n-1}$; the
merged collection covers $G$, properly covers every fibre, and is a chain in Bing and Jones's sense,
with nonadjacent links at positive distance (checked by hand; the certificate uses the library's
notion of chain, in which nonadjacent links are disjoint).
Take $U$ a small neighbourhood of $\overline{g_c\cap e_n}$ inside $e_{n-1}$. Then $g_c-U$ misses
$e_n$, and the conclusion fails at $c$ for every separation, although no fibre lies in $U$.
Certificates: \lean{bj_theorem5_printed_false_general} (from every instance with three or more
links) and \lean{bj_theorem5_printed_false} (a fibre inside $U$). The printed proof passes from
separations of $G-d_j$ to separations of $G-U$.

\fld{Consequence} The theorem's one application (p.~185, in Step~2(c) of Theorem~10: ``By Theorem 5
there is an open set $U_i$ \ldots'') is existential, and a small enough $U_i$ exists. Theorem~10 is
unaffected. The repaired form asks that $f$ be open and that no fibre meet a link only inside
$\bar U$; it is verified (\lean{bj_theorem5}, \lean{theorem5_repaired}).
\end{defect}

%----------------------------------------------------------------------
\begin{defect}{Bing 1948, Theorem~8, p.~737}{False as printed. Counted (one false statement).
Theorem~9's proof uses it. Ledger A.1.3.}{1}{0}\label{it:b48-8}

\fld{Printed} ``\textsc{Theorem 8.} \emph{Suppose the chain $D_n$ ($n=1,2$) is crooked in the chain
$E$, each end link of $E$ contains a link of $D_n$, and no proper consolidation of $D_n$ has these
properties. Then $D_1$ follows a pattern in $E$ that $D_2$ follows in $E$, and the chain whose
links are the sums of corresponding links of $D_1$ and $D_2$ is crooked in $E$.}'' The proof:
``That the chain whose links are the sum of the corresponding links $D_1$ and $D_2$ is crooked in
$E$ may be proved by induction on the number of links in $E$.''

\fld{What fails} The unions of corresponding links need not form a chain. Certificate
(\lean{not_theorem8_literal}): in a finite discrete subspace of $\mathbb N$, two chains of six links
(the fewest possible, \lean{theorem7_minimal}) are crooked in a four-link chain $E$, have their
first links in $e_1$ and last links in $e_4$, and follow Bing's pattern in $E$ in the same
direction; the second and fourth unions meet (at the point $20$), so the unions form no chain, and
pairing the links in opposite directions gives unions that do not refine $E$. Bing's chains have
open links that need not be connected (p.~729: ``We have not supposed that the links of a chain
are necessarily connected''), so the discrete example applies; minimality holds because a proper
consolidation would have fewer links. The mechanism is that the pattern revisits levels, so links
of $D_1$ and $D_2$ at non-adjacent positions can meet.

\fld{Consequence} The proof of Theorem~9 (p.~737) begins ``By Theorems 7 and 8 there are chains
$G_1$ and $G_2$ \ldots\ and the chain $H$ \ldots\ whose links are the sums of corresponding links
of $G_1$ and $G_2$ is crooked in $E$.'' The corrected Theorem~8 (\lean{Fold.theorem8}, with the
orientation made explicit and the chain property of the unions as a hypothesis) suffices.
Theorem~9 is verified with its mesh clause weakened: each link of $F$ lies in the sum of three
consecutive links of $E$, where Bing has ``the sum of two adjacent links of $E$''
(\lean{Fold.theorem9}). The printed two-link form is not formalized.
\end{defect}

%----------------------------------------------------------------------
\begin{defect}{Moore 1919, Theorems C and D}{Theorem C false as printed; the proof of Theorem D
defective, statement true. Counted (one false statement, one defective proof). Ledger
A.15.}{1}{1}\label{it:m19}

\fld{Theorem C} It concerns rulings of a closed curve with respect to a quadrilateral $ABCDA$. A
closed curve that touches a side from within meets $ABCDA$ in one point and admits no ruling of the
required kind. Moore offers an ``Indication of proof \ldots\ given instead of a tedious formal proof
covering all possible cases'', which misses this case. Certificate:
\lean{not_m19_theoremC_printed}, in every model of Moore's Axioms~1--8. Three corrected forms are
verified: absorbing finitely many disjoint cross-cuts (\lean{m19_theoremC_crosscuts}), a curve
within $ABCDA$ (\lean{m19_theoremC_within}), and the printed statement with touching from within
excluded (\lean{m19_theoremC_corrected}). Touching from without is harmless.

\fld{Theorem D} Moore's proof applies Theorem~C to a curve $\bar J$ that $g_1$ may touch, so it
rests on the false case of C; its induction over $g_2,\ldots$ needs D itself; and the construction
of Fig.~2 is unargued (``It easily follows''). Verified (\lean{m19_theoremD}) by an induction over a
finite subcover: at each step a curve is built about the uncovered part, crossing $ABCDA$ at every
common point, and the corrected Theorem~C is applied. Whether Moore's $\bar J$ could be re-chosen
so that it is crossed and not touched is not decided.

\fld{Consequence} The paper's conclusion, that a space satisfying Moore's Axioms~1--8 is
homeomorphic to the plane, is verified (\lean{m19_plane}), with Theorems E, G and the final
paragraph reached by a route different from Moore's (\S\ref{sec:unargued}).
\end{defect}

%======================================================================
\section{Other proofs that do not go through}\label{sec:other}

The items of this section concern supporting results of plane topology used inside the proofs of
\S\ref{sec:path}. In each, the printed argument fails in a case the statement covers, and the
result is proved by a different argument.

%----------------------------------------------------------------------
\begin{defect}{Whyburn 1939, Theorem~14, p.~747}{Proof defective; statement true. Counted (one
defective proof). Ledger A.21.}{0}{1}\label{it:w39}

\fld{Printed} ``Clearly we may suppose $R$ bounded, since an inversion of the plane can be used to
throw $R$ into a bounded region complementary to the image of $M$.''

\fld{What fails} The centre of the inversion must lie in a complementary domain of $M$ other than
$R$. When $M$ has no bounded complementary domain, for instance when $M$ is a closed disk or any
continuum that does not separate the plane, there is none. This is a natural case of the theorem,
not a degenerate one.

\fld{Formalization} Verified for bounded and unbounded domains alike, without inversion
(\lean{whyburn39_theorem14}). Moore's curve $C$ is replaced by a joining lemma from Janiszewski's
theorem, and Whyburn's curve $K$ and limit continuum $H$ by two facts: a point where infinitely
many disjoint connected sets accumulate has no finite order, and Whyburn's (3.2).
\end{defect}

%----------------------------------------------------------------------
\begin{defect}{Moore 1925, Theorem~21: Axioms 6 and 7}{Proof defective; statement true. Counted (one
defective proof). Ledger A.16.}{0}{1}\label{it:m25-21}

\fld{Printed} ``With the aid of results established in my paper \emph{Concerning a set of
postulates for plane analysis situs} \ldots\ it easily follows that Axiom 6 and 7 also hold true.''

\fld{What fails} Moore 1919 proves nothing for the axiom system $\Sigma_2$ beyond the assertion of
its introduction. The step therefore rests on the equivalence of $\Sigma_1$ and $\Sigma_2$ that
Moore 1916 asserts on p.~163 without proof (``It is easy however to show that though $\Sigma_1$
and $\Sigma_2$ are not absolutely equivalent, they are equivalent with respect to point and limit
point''). The formalization neither proves nor needs that equivalence.

\fld{Formalization} Verified (\lean{m25_theorem21}): Axioms 6 and 7 are proved directly, with
cross-cuts and Moore 1925 Theorems 10, 12, 15, 18 and 19. The other three ``easily'' claims of the
proof are confirmed. Moore 1925 Theorem~22, the plane quotient used for Theorem~A, is verified from
Moore 1919 for $\Sigma_1$ only (\lean{m25_theorem22}).
\end{defect}


%======================================================================
\section{Citations that do not carry what they are cited for}\label{sec:cite}

None of these is counted. The first bears on the classification (Item~\ref{it:unique}); the others
were avoided or repaired.

%----------------------------------------------------------------------
\begin{defect}{Lewis 1983 and 1996: uniqueness attributed to work that leaves it open}{Citation:
the cited work leaves the claim open. Ledger A.27.}{0}{0}\label{it:lewis}
Lewis \cite[p.~80]{Lw83} poses the uniqueness question as ``Question [28]'', where [28] is Lewis
1985 (then submitted), and asserts without proof that it ``has an affirmative answer if all of the
decomposition elements in each of $G_1$ and $G_2$ are terminal continua''. Lewis \cite[p.~173]{Lw96} asserts the uniqueness of the continuum ``obtained either by
shrinking out maximal terminal pseudo-arcs or by enlarging each point to a terminal pseudo-arc'',
on the strength of an observation of his own and of Lewis 1985, Theorem~2 [24], which gives
existence and lifting only. Lewis 1985 \cite[p.~93]{Lw85}: ``We still do not know in general
whether $\hat M$ is unique''; and \cite[p.~99, Question 2]{Lw85}: ``If $\hat M_1$ and $\hat M_2$
are two continuous curves of pseudo-arcs with homeomorphic decomposition spaces, are $\hat M_1$
and $\hat M_2$ homeomorphic?'' The case the classification needs is Domokos--Prajs Theorem~6.1
(Item~\ref{it:unique}).
\end{defect}

%----------------------------------------------------------------------
\begin{defect}{Domokos--Prajs 2022, Example 6.1: terminality of Bing and Jones's
decomposition}{Citation: the cited work does not prove the claim; not used. Ledger A.27.}{0}{0}\label{it:dpex}
p.~13: Bing and Jones ``constructed the \emph{circle of pseudo-arcs}, denoted here by $\Psi$, which
is a homogenous, decomposable and nowhere locally-connected continuum in the plane admitting a
continuous decomposition into terminal pseudo-arcs.'' Bing and Jones do not prove that the
elements are terminal. This supports ``It is known'' in the proof of the ``only if'' direction of
Theorem~6.1 (p.~14). The formalization proves that direction differently: a circle of pseudo-arcs
is homogeneous (Bing--Jones Theorem~10), so Jones's decomposition gives an atomic map
(\lean{jones55_c}), and the property passes to every homeomorphic space
(\lean{dp22_theorem61_iff'}).
\end{defect}

%----------------------------------------------------------------------
\begin{defect}{Hoehn--Oversteegen 2016, pp.~191 and 211: a research announcement cited for a
proof}{Citation: an announcement without proofs; avoided. Ledger A.26.}{0}{0}\label{it:lelek}
p.~191: ``Because $X$ has span zero, it is tree-like [31]''; p.~211: ``Since $X$ has span zero, by
[31] we have that $X$ is tree-like.'' Here [31] is Lelek 1979 \cite{Le79}, a research announcement
(p.~631: ``The purpose of the present paper is to announce some results \ldots\ A complete version
will be published elsewhere''). Its Propositions~1--4 are called ``rather straightforward'' and
are not proved; its Lemma~5 and Theorem~6 carry no proof; its Corollary~9 drops
Theorem~6's hypothesis ``$\dim X\le1$''; its Corollary~10, that span zero implies tree-like, is
obtained in one paragraph through three citations. On p.~212 Hoehn and Oversteegen cite a published
proof of a stronger statement ([46], Oversteegen--Tymchatyn 1984). The formalization does not need
tree-likeness: for plane continua a grid graph serves, since hereditary indecomposability gives
empty interior and radial projection in the squares of a fine grid maps the continuum into the
union of their edges.
\end{defect}

%----------------------------------------------------------------------
\begin{defect}{Fugate 1966: Miller's and Sorgenfrey's theorems cited in over-general
form}{Citation: over-general as cited; harmless in context. Ledger A.4.3, A.4.4.}{0}{0}\label{it:fugate-cite}
As cited, Miller [6, Theorem~1.5, p.~180] says: if $A\cup B$ contains a non-unicoherent subcontinuum
$D$, then $D\subseteq A$ or $D\subseteq B$. This is false: the closed disk is the union of two
half-disks, and its boundary circle lies in neither (\lean{miller_1_5_false}). What Fugate needs is
hereditary unicoherence of $A\cup B$ for chainable $A$, $B$ with connected intersection, proved
directly (\lean{isHereditarilyUnicoherent_union}). Miller's own statement has not been checked
against the citation. As cited, Sorgenfrey [8, p.~440] says that three intersecting continua, none
contained in the union of the other two, have a union that is a triod; a circle is the union of
three arcs. Inside a hereditarily unicoherent continuum, which is Fugate's situation, it holds.
In Fugate's setting the continua are also a-triodic, and neither citation is needed in the
generality printed.
\end{defect}

%----------------------------------------------------------------------
\begin{defect}{Hagopian 1975, 1976 and 1984: citations that do not cover the step}{Citation:
the cited result does not cover the step; each gap filled. Ledger A.22--A.24.}{0}{0}\label{it:hag-cite}
\begin{itemize}
\item Hagopian 1976, p.~344: ``Since $M$ is indecomposable and homogeneous, no proper subcontinuum
  of $M$ separates $R^2$ [5, Lemma 1 (second paragraph of proof)]''. The cited argument carries a
  \emph{non-unicoherent} continuum to every composant; for a separating one it needs Borsuk's
  theorem that plane separation is a topological invariant, which neither paper states. Verified
  without Borsuk (\lean{no_proper_subcontinuum_separates}, through
  \lean{exists_separating_subcontinuum_of_homeomorph}).
\item Hagopian 1984, p.~346: ``since the elements of $\mathfrak D$ are tree-like, if $M$ is planar,
  then $M/\mathfrak D$ is planar [21]''. Moore 1925 needs elements that do not separate the plane;
  that tree-like continua do not is not stated. Filled in the plane by the same theorem.
\item Hagopian 1975, Lemma~2 (pp.~36--37): ``[10, Theorem 3, p.~216]'' states that Kuratowski's
  family has the order type of the interval, not that the continuum is of type A$'$. Filled by a
  countable dense subfamily, Cantor's theorem and the map $k(y)=\inf\{q:y\in D_q\}$. Hagopian 1984
  p.~335 repeats the citation with the same gap. Both also cite Thomas 1966 (Theorem~10, p.~15),
  which was not in hand and is not checked; the gap is recorded for the Kuratowski citation only.
\item Hagopian 1975, paragraph 13: ``[9, Theorem 2]'': Jones 1955 Theorem~2 implies, but does not
  state, that a unicoherent homogeneous plane continuum is indecomposable. Supplied.
\item Hagopian 1984, Lemma~9: the citation (1975, paragraphs 5--12) treats interior layers; the
  lemma claims every element, end sets included. True through an extension, not stated.
\end{itemize}
\end{defect}

%----------------------------------------------------------------------
\begin{defect}{Other citations}{Citation; harmless. Ledger A.11, A.18, A.21, A.24.}{0}{0}\label{it:cite-misc}
\begin{itemize}
\item Jones 1955, Theorem~2, last step: ``Hence by [5]'': the result used is [3] (Jones 1949), not
  [5] (Jones 1951).
\item Jones 1949, reference [8]: ``Theorem 9.1, p.~61'' is Whyburn's (9.2); (9.1) is the lemma
  before it. Reference [4]: Menger's theorem characterizes the circle by \emph{all} points being
  of order 2, and Jones applies it to \emph{a} point of order 2; homogeneity and the invariance of
  order are tacit (supplied, \lean{jones49_order_two}).
\item Whyburn 1939, p.~748, footnote 14: ``by a theorem of R. L. Moore'' gives no theorem number. In
  substance it is Moore 1916 Theorem~43 in corrected form (Item~\ref{it:m16-degen}).
\item Hagopian 1984, p.~346: ``In 1968 F. B. Jones [16]'', where the reference list gives 1969.
\end{itemize}
\end{defect}

%======================================================================
\section{Proof steps that fail locally}\label{sec:local}

In each counted item of this section a step of the printed proof is false as written, but the
argument stands for the main case: a short argument inside the same proof repairs the step, or
(Item~\ref{it:bj7proof}) the failure is confined to chains of at most three links, which need
another construction.

%----------------------------------------------------------------------
\begin{defect}{Bing--Jones 1959, Theorem~7, proof, p.~177}{Proof defective; statement true. Counted
(one defective proof). Ledger A.25 (third reading).}{0}{1}\label{it:bj7proof}
The reduction to the chain $D'=U_0,U_1,U_2,d_3,\cdots,d_{n-2},U_{n-1},U_n,U_{n+1}$ works for $n\ge4$.
For $n=3$ its conditions cannot all hold: a link serving as both $U_2$ and $U_{n-1}$ contains
$P\cdot d_2$ and meets $U_4$ (certificate \lean{bj_theorem7_proof_n3}). Chains of one or two links
are not treated. Unargued as well: ``We suppose with no loss of generality that $r<s$''; the
special case $A\subset e_1$; the existence of the crooked sequence $D_1=D,D_2,\cdots$; and the
non-emptiness of $P\cdot(d_1-d_2)$. Verified by another construction, following the pattern along
a walk and folding, with no special case (\lean{bj_theorem7}); the second sentence holds for
$n\ge3$ with a collar chain of $n+4$ links in place of $D'$.
\end{defect}

\begin{defect}{Hoehn--Oversteegen 2016, Proposition~19, p.~208}{Proof defective; statement true.
Counted (one defective proof). Ledger A.26 (fourth reading).}{0}{1}\label{it:ho19}
``If $D\subset F_1\cup F_2$, then \ldots\ $\varphi(D)$ is contained in a component $\widetilde D$ of
$G\setminus\pi(S_{i_0})$. Moreover, $\varphi(D)\subset\pi(P_1\cup P_2)$, so
$\partial\widetilde D\cap\pi(\alpha_{i_0})=\emptyset$ by (S6$'$) for $S$. This means that either
$\partial\widetilde D\subset\pi(\beta_{i_0})$ or $\partial\widetilde D\subset\pi(\gamma_{i_0})$.'' The
middle sentence is false. Certificate \lean{ho16_p208_false}, verified on the graph $G$: a broken
stairwell structure of height 3 with its pit at level 2 on $G=[0,10]$, every sheet at height
$\tfrac12$, with $\pi(S_1)=[4,10]$, $\pi(S_2)=[4,6]$, $\pi(S_3)=G$, $\pi(\alpha_2)=\{4\}$,
$\pi(\gamma_2)=\{6\}$, $\pi(P_1)=[1,2]\cup[\tfrac{11}2,10]$, $\pi(P_2)=[1,2]\cup[\tfrac{11}2,6]$;
the point $x=\tfrac32\in\pi(P_2)\setminus\pi(S_2)$ lies in the component $[0,4)$ of
$G\setminus\pi(S_2)$, whose frontier point $4$ is in $\pi(\alpha_2)$. That the copies of $[1,2]$ in
$F_1$ and $F_2$ form a component $D\subseteq F_1\cup F_2$ of the fold with
$\widetilde D=[0,4)$ is a hand argument; the fold is not constructed in the certificate. The
conclusion of the step holds, and the sentence fails only for components that the restriction on
p.~209 discards. Verified: \lean{ho16_prop19}, \lean{ho16_prop19_full}, with a case split.
\end{defect}

\begin{defect}{Moore 1916, Theorem~10}{Proof defective; statement true. Counted. Ledger
A.12.}{0}{1}\label{it:m16-10}
``$R_1R_2\cdots R_kR$ is a simple chain'' need not hold. Repaired by taking a shortest sub-chain
(\lean{fa_theorem10}).
\end{defect}

\begin{defect}{Moore 1916, Theorem~15, part III}{Proof defective; statement true. Counted. Ledger
A.12.}{0}{1}\label{it:m16-15}
``since $y$ precedes every link of $C_n$ that contains $P$ it follows that $P_B$ precedes $P$''
fails at stage $n$, because $P_B$ can lie in the next link, which may contain $P$. Repaired by
choosing $n$ large; the two ``it is clear'' steps of part III are supplied
(\lean{fa_theorem15}).
\end{defect}

\begin{defect}{Moore 1916, Theorem~28, last step}{Proof defective; statement true. Counted. Ledger
A.13.}{0}{1}\label{it:m16-28}
The last step: ``the point-set $A_2C+CP+PD+DB_2$ contains as a subset an arc $AXB$ satisfying
conditions (1), (2), and (3)''. The arc $CA_2$ runs in $R_A$, which straddles $J$, so its part
before $J$ need not lie in $M$. Repaired by taking $A$ as the first point of $CA_2$ outside $M$,
which lies on $J$; the point $P$ is not needed (\lean{fa_theorem28}).
\end{defect}

\begin{defect}{Moore 1916, Theorems~33--34, Kline's proof, p.~150}{Proof defective; statements true.
Counted once. Ledger A.13.}{0}{1}\label{it:m16-33}
Two proofs are printed; Kline's general proof has a gap: the first point of $DC$ on $CAB$ may be
$B$. Filled (\lean{fa_theorem33}, \lean{fa_theorem34}). Theorem~34 is verified for any region.
\end{defect}

\begin{defect}{Moore 1916, Theorem~36}{Proof defective; statement true. Counted. Ledger
A.14.}{0}{1}\label{it:m16-36}
The arc $BC$ can leave $R$; it is cut at its first point outside $R$ (\lean{fa_theorem36}). Moore
1925 Theorem~19 needs, in addition, that the interior lies in $R$
(\lean{fa_theorem36_interior_subset}).
\end{defect}

\begin{defect}{Moore 1925, Theorem~11, pp.~421--422}{Proof defective; statement true under the
corrected definition (Item~\ref{it:m25-def}). Counted. Ledger A.16.}{0}{1}\label{it:m25-11}
``If $B$ is not identical with $J$ then it is a proper subset of $J$. It easily follows, with the
help of Theorems 10 and 6, that $B$ is a simple continuous arc''. Theorem~6 needs $B$ connected,
which is not shown. Repair: $B$ lies in a sub-arc $A$ of $J$ missing an element of $J-B$, and $A$
would separate $G$, against Theorem~9 (\lean{frontier_eq_of_scc}, \lean{m25_theorem11}).
\end{defect}

\begin{defect}{Moore 1922, Theorem~4, necessity, last step, p.~237}{Proof defective; statement
true. Counted. Ledger A.20.}{0}{1}\label{it:m22}
The points $Y_m$, $Y_n$ lie in $R$, not in $M$, so ``there exists no connected subset of $M$ of
diameter less than $\varepsilon$ that contains both $Y_m$ and $Y_n$'' proves nothing. Repaired
within the proof: take points of $\bar R_n$ on the middle square and use $\varepsilon/2$
(\lean{moore22_theorem4}).
\end{defect}

\begin{defect}{Whyburn, \emph{Analytic Topology}, (9.1)}{Proof defective; statement true. Counted.
Ledger A.19.}{0}{1}\label{it:w42}
The proof applies [I, (10.2)], which requires $N\cdot G$ compact, to $U\cdot(Q\cdot R_a+p)$, which
need not be compact. The case $N\cdot G=N$ is not treated, and ``$Q\cdot R_a+p$ is a generalized
continuum'' is asserted. The repair runs boundary bumping in a compact neighbourhood; the openness
of $R$ is not needed (\lean{whyburn_9_1}).
\end{defect}

\subsection*{Local defects not counted}
These are defective constructions or reasons inside proofs whose step is nonetheless right, or is
repaired in one line.
\begin{itemize}
\item Bing--Jones 1959, Theorem~3 (p.~175): even with $F(w,m)$ read for the printed $F(v,m)$, the
  chain $D$ is not a chain as written: $A\cdot e_{j-1}$ can meet $f_{v+1}$ and $B\cdot e_{r+1}$ can
  meet $f_{w-1}$, which are not adjacent in $D$. Repaired by merging (\lean{bj_theorem3}).
\item Bing--Jones 1959, Theorem~8 (p.~178): the printed $e_k$ need not form a chain, since
  $e(a_i)_k$ is controlled only over $I_{a_i}$ (argued, not certified). Repaired by restricting each
  piece to the fibres over $I_{a_i}$ (\lean{bj_theorem8}).
\item Bing--Jones 1959, Theorem~9 (pp.~179--180): the condition on $\varepsilon$ concerns subsets
  of $Q$, while the links of $F$ are not subsets of $Q$; it secures neither the disjointness of
  links in different middles nor that every link falls under a rule. Repaired with links whose
  closures lie in links of $E$ and a smaller mesh (\lean{bj_theorem9}).
\item Hoehn--Oversteegen 2016, Proposition~10, second part (p.~190): ``If $C'\neq C$, then there
  must be a point $x\in\partial C'\cap C$. But since $A'$ has consistent complement relative to $B$,
  we must have $x\in B$''. Consistency gives $\partial C'\subseteq B$ only if $\partial C'$ meets
  $B$, which is not shown. Repaired by shrinking $V$; the formal proof avoids the step
  (\lean{ho16_prop10}).
\item Hoehn--Oversteegen 2016, Theorem~15, Claim~15.1, property (4) (p.~195): ``because $x$ is not an endpoint
  of $G$ by (5), and $\pi(M'\cap V)$ is open for some neighborhood $V$ of $(x,y)$ by (3)''. Here
  $(x,y)\in T$, and (3) speaks only of points of $M'\setminus T$; the openness comes from
  irreducible separation. The step is true. Also p.~198, (S4): an incomplete case split (a
  component around a cycle of $G$ is in neither case; the same chain argument covers it); and
  p.~197: nothing excludes $b\in J_j$. The formal proof of Theorem~15 avoids all of pp.~194--198
  (\lean{ho16_theorem15}).
\item Hagopian 1976, Lemma~4, first paragraph (p.~341): ``Since the decompositions of $E$ and $F$
  are minimal, each element of the decomposition of $E$ (except possibly $k^{-1}(1)$) is an element
  of the decomposition of $F$.'' Minimality gives it only for layers in $d^{-1}(u)$ with $u<s$
  (\lean{hagopian76_lemma4_layers'}); a plane model without homogeneity has a layer swallowed.
  Verified in a strong form through Lemma~5 (\lean{hagopian76_far_layer}).
\item Hagopian 1976, $\varepsilon_1$ (pp.~345--346): ``the chain formed by taking
  $\varepsilon_1$-neighborhoods \ldots\ follows $h[E]$'' does not follow from the four printed
  families; true after shrinking $\varepsilon_1$. Likewise property (D) of T9.
\item Hagopian 1984, Lemma~15 (p.~341): ``By Lemma 3'' is insufficient as printed (the layer through
  a point of $h[Z]$ need not lie in $\Delta(\mathcal P)$); verified with a smaller $\varepsilon$
  (\lean{hagopian84_lemma15}). The cover (p.~343): ``it follows that $\bigcup\mathfrak B$ is a closed
  open subset of $M$'' is not an argument over uncountably many $Z$; a finite subcover is proved
  directly (\lean{hagopian84_cover}, \lean{sUnion_cover84_eq_univ}). Disjointness, first case:
  ``by (2), (3), and Lemma 7, $A=B$'' leaves the case $S_A\subseteq Y$ open; supplied by a transfer
  lemma.
\item Oversteegen--Tymchatyn 1982, Corollary~3.5 (pp.~139--140): ``there is a continuum $B\subset
  X$ such that $[0,5/8]\subset\pi_1(h_K(B))\subset[0,3/4)$'': the first inclusion is impossible as
  printed, since $h_K(B)\subset K\subset[\varepsilon,3/4]\times R$. Of the two readings only
  $[0,5/8]\subset\pi_1(B)$ gives what Theorem~3.2 needs; the existence of $B$ is unargued and is
  supplied through hereditary indecomposability. ``Without loss of generality \ldots\
  $\pi_1(A)=[0,1/2]$ and $\pi_1(X)=[0,1]$'' is not a rescaling; supplied by a shear and a
  homeomorphism.
\item Jones 1975, (d), second sentence (p.~58): ``$c'$ has the properties possessed by $c$'' holds
  only if the Effros homeomorphism preserves the direction of $T$; adding $d(a,ce)$ to the minimum
  repairs it. The triod paragraph does not yield the contradiction as printed.
\item Fugate 1969, Lemma~3: Lemma~2 gives $K\cup R$ irreducible between some $q\in K$ and some $p$;
  getting $p\in R\setminus K$ needs $R\not\subseteq K$, which is not argued. Supplied from
  Sorgenfrey's theorem.
\end{itemize}

%======================================================================
\section{Statements false because of their wording}\label{sec:wording}

In each item a definition or hypothesis says less than the author meant, and a configuration the
author plainly did not have in mind refutes the printed statement. The intended statement is
verified in each case.

\begin{defect}{Moore 1925, the definition of complementary domain (p.~417) and Theorem~2}{False as
printed. Counted (one false statement). Ledger A.16.}{1}{0}\label{it:m25-2}
``A domain of elements of $G$ is said to be a complementary domain of a closed set $H$ \ldots\
provided the boundary of $D$ is a subset of $H$.'' Neither disjointness from $H$ nor maximality is
required, so $G$ is a complementary domain of every closed set, and ``the unbounded complementary
domain'' is not unique. Theorem~2 is false under this definition
(\lean{not_m25_theorem2_printed}): take the decomposition into points, $M$ a closed disk, $D$ the
open disk minus a radius, and $E$ the plane minus the radius; the outer boundary is the radius,
which is cut by each of its interior points. Read as components of the complement, Theorem~2 holds
(\lean{m25_theorem2}), and the unbounded one is unique (\lean{m25_unbounded_unique}).
\end{defect}

\begin{defect}{Moore 1925, the definition of simple closed curve (p.~417) and Theorems~10 and
11}{False as printed. Counted (two false statements). Ledger A.16.}{2}{0}\label{it:m25-def}
``A bounded subset $H$ of the set $G$ is said to be a simple closed curve \ldots\ provided it is
disconnected by the omission of any two of its elements.'' Neither closedness nor connectedness is
required, nor more than one element. Two disjoint segments (decomposition into points) qualify; for
them Theorems~10 and 11 fail: they are not a union of two arcs, and their complement is connected
(\lean{not_m25_theorem10_printed}, \lean{not_m25_theorem11_printed}). With ``closed, connected,
nontrivial'' added, both are verified (\lean{m25_theorem10}, \lean{m25_theorem11}).
\end{defect}

\begin{defect}{Moore 1925, Theorem~7, p.~420}{False as printed, by a misprint. Counted (one false
statement). Ledger A.16.}{1}{0}\label{it:m25-7}
``then $p$ is a limit element of the set $H$'' should read $x$. The printed form is refuted
(\lean{not_m25_theorem7_printed}); the corrected form is verified (\lean{m25_theorem7}).
\end{defect}

\begin{defect}{Mullikin 1922, Theorem~5, p.~160}{False as printed. Counted (one false statement).
Ledger A.17.}{1}{0}\label{it:mu5}
``\ldots\ such that $M_1$ and $M_2$ have in common only $K_1$ and $K_2$, where $K_1$ and $K_2$ are
mutually exclusive connected sets, then $S-M_1-M_2$ is the sum of just two mutually exclusive,
connected domains.'' Mutually exclusive connected sets need not be separated: take $M_1=M_2$ a
segment, $K_1$ a half-open half and $K_2$ the closed other half; their union is connected, and the
complement of the segment is connected (\lean{not_mullikin_theorem5_printed}). The intended
hypothesis is that $K_1$, $K_2$ be closed; Moore's use ($K_1=\bar p$, $K_2=\bar q$) satisfies it.
Verified in that form (\lean{mullikin_theorem5}), by logarithms and Borsuk's theorem instead of
polygons.
\end{defect}

\subsection*{Omitted hypotheses and definitions, not counted}
\begin{itemize}
\item Hoehn--Oversteegen 2016, pp.~193--194: ``if $k$ is odd then $\pi(S)=G$'' and ``$\pi$ maps $S$
  homeomorphically onto $G$'' omit ``$G$ connected''; with the paper's convention (``$G$ will
  denote a (not necessarily connected) graph'', p.~184) both fail for $G=[0,1]\cup[2,3]$,
  $S=[0,1]\times\{\tfrac12\}$ (\lean{ho16_p193_remark_false}).
\item Hoehn--Oversteegen 2016, Proposition~17 (p.~200): the hypothesis $k\ge3$ is not printed (for
  $k\le2$ the conclusion has height $\le0$); harmless, since Theorem~20 applies it to odd $k\ne1$.
\item Moore 1916, Theorem~31: the order $C\,E\,B\,G\,D$ needs $B\ne C,D$. Theorem~32: nothing keeps
  $C\in R$ off the arc $AB$; proved under $R\cap AB=\emptyset$. ``The interior of $J$'' is used from
  Theorem~23 on and never defined; supplied as the region bounded by $J$.
\item Moore 1919: the double-ruling definition (p.~172, footnote) requires equal numbers of arcs in
  the two families, which Moore's own Fig.~1 (9 and 8) violates; padding shows the requirement costs
  nothing. It omits the disjointness of the $Y$-arcs; the ``subdivisions'' are never defined.
\item Moore 1925: the definition of boundary (p.~417) is empty in its second clause as printed;
  read ``of $G$'', it is the frontier. Theorem~3's statement does not say that the region contains
  $p$, which the proof gives and Theorem~4 needs.
\item Moore 1920, Definition~4 (p.~342) admits a single point; see Item~\ref{it:m20}.
\item Hagopian 1976, p.~340: the end-set property ``[14, Theorem 8, p.~14]'' is true for the
  decomposition and false for an arbitrary admissible map.
\item Mullikin 1922, Theorem~4 (p.~154): true for disjoint $M_1$, $M_2$ under her convention that the
  empty set is connected (\lean{mullikin_theorem4_disjoint}), false under Mathlib's
  (\lean{not_mullikin_theorem4_disjoint_mathlib}).
\end{itemize}

%======================================================================
\section{Statements false at degenerate cases}\label{sec:degenerate}

Each statement in this section fails only at an empty set, a single point, an endpoint, or a
chain of at most three links, and holds once that case is excluded. In most of them the proof
begins by choosing a point that the hypotheses do not supply. Bing 1948 Theorem~4 is different: its
proof claims the small case as easy (Item~\ref{it:b48-4}). Moore 1916 Theorem~43 has no printed
proof.

\begin{defect}{Jones 1948, Theorems~6 and 10}{False as printed. Counted (two false statements).
Ledger A.10.}{2}{0}\label{it:j48}
Theorem~10 (p.~407): ``\emph{In order that a compact continuum $M$ be indecomposable it is
necessary and sufficient that $M$ contain an open subset $D$ of $M$ such that if $x$ is a point of
$D$ and $y$ is a point of $M-x$, then $M$ is not aposyndetic at $x$ with respect to $y$.}''
$D=\emptyset$ satisfies the condition vacuously in every continuum, so the statement says every continuum is indecomposable; an arc refutes it
(\lean{theorem10_printed_iff}, \lean{not_theorem10_printed_unitInterval}). Theorem~6 has the same
defect: with $D=M_0=\emptyset$ the hypotheses hold vacuously and the conclusion asks for a point of
$D$ (\lean{not_theorem6_printed}). Both are verified with $D$ nonempty. Theorems~11, 12, 13 and 19
carry the same tacit hypothesis and are not listed as false, because the intended reading is plain.
\end{defect}

\begin{defect}{Moore 1916, Theorems~4, 14, 19, 23, 35 and 43, and the limit-point criterion of
p.~139}{False as printed. Counted (seven false statements). Ledger A.12--A.14.}{7}{0}\label{it:m16-degen}
\begin{itemize}
\item Theorems~4, 19, 23 and 35: adjoining $\emptyset$ to the regions preserves Axioms~1--8 and
  every derived notion, and each statement fails for the empty region. Moore's proofs begin by
  choosing a point of $R$. Verified for $R$ nonempty (\lean{fa_theorem4}, \lean{fa_theorem19},
  \lean{fa_theorem23}, \lean{fa_theorem35}).
\item Theorem~14: the proof opens ``For each $n$, $M_n$ contains at least one point'', which is not
  a hypothesis; $M_n=\emptyset$ refutes it (\lean{not_fa_theorem14_printed}); verified with $M_n$
  nonempty (\lean{fa_theorem14}).
\item Theorem~43: false for $G=\emptyset$ in every model (\lean{not_fa_theorem43_printed}). It has no
  printed proof; the corrected form is verified through a dual of Theorem~41 for intersections of
  interiors (\lean{fa_theorem43}).
\item The limit-point criterion on an arc (p.~139, unnumbered) fails at an endpoint, where it holds
  vacuously (\lean{not_fa_sec5_limitPt_printed}). Corrected and verified.
\end{itemize}
\end{defect}

\begin{defect}{Moore 1920, Theorem~4}{False as printed. Counted (one false statement). Ledger
A.16.}{1}{0}\label{it:m20}
Definition~4 (p.~342) of a simple closed curve has ``closed, connected, and bounded'' but admits a
single point, for which Theorem~4 is false (\lean{not_moore1920_theorem4_singleton}); so is Moore
1925 Theorem~11 (\lean{not_m25_theorem11_singleton}).
\end{defect}

\begin{defect}{Jones 1949, Theorems~1 and 2}{False as printed. Counted (two false statements).
Ledger A.18.}{2}{0}\label{it:j49}
Theorem~1: ``If a compact, plane continuum $M$ is both homogeneous and aposyndetic, then $M$ is a
simple closed curve'' fails for a one-point continuum (\lean{not_jones49_theorem1_point}). Jones's
footnote 3 refers the reader for definitions to Moore's book or Whyburn's; Whyburn's convention
(\emph{Analytic Topology}, p.~15: ``A compact connected set will be called a continuum'') admits
one-point continua. Theorem~2 fails likewise: a point is
homogeneous and has no weak cut point (\lean{not_jones49_theorem2_point}). Both are verified for
nondegenerate continua (\lean{jones49_theorem1}, \lean{jones49_theorem2}).
\end{defect}

\begin{defect}{Hoehn--Oversteegen 2016, Theorem~1$'$, p.~212}{False as printed. Counted (one false
statement). Ledger A.26.}{1}{0}\label{it:ho1p}
Theorem~1$'$ (``surjective semispan zero'' in place of ``span zero'') omits ``non-degenerate'' and
fails for a one-point continuum (\lean{ho16_theorem1prime_printed_false}). It is stated correctly
nowhere, although p.~212 presents it as what ``our proof of Theorem 1 in fact establishes''.
Theorem~1 as restated on p.~211 omits ``non-degenerate'' too
(\lean{ho16_theorem1_printed_false}); it is recorded as a misprint, because p.~180 states it
correctly.
\end{defect}

\begin{defect}{Bing 1948, Theorem~4, at $m=3$, p.~733}{False as printed for $m=3$; true for $m\ge4$.
Counted (one false statement). Ledger A.1.2.}{1}{0}\label{it:b48-4}
``\textsc{Theorem 4.} \emph{If $D=[d_1,d_2,\cdots,d_n]$ is a chain crooked in the chain
$E=[e_1,e_2,\cdots,e_m]$ and $D(r,s)$ is a subchain of $D$ such that a link of $D(r,s)$ intersects
$e_1$ and a link of $D(r,s)$ intersects $e_m$, then there is a chain $F$ such that $F$ is a
consolidation of $D$, each element of $F$ is a subset of the sum of two adjacent links of $E$,
$d_r$ is contained in only the first link of $F$, and $d_s$ is contained in only the last.}'' The
proof: ``Since the case of our theorem where $m\le4$ is easily established''. Two failures at
$m=3$. With $r=s$: three intervals in $\mathbb R$ crooked in three intervals, the middle fine link
meeting the first and last coarse links; $d_r$ in both the first and the last link of $F$ forces
$F$ to have one link, all of $D$, which lies in no two adjacent links of $E$
(\lean{theorem4_two_degenerate_false}, \lean{not_theorem4Literal}). With $r<s$: in $\mathbb N$
with the discrete topology, coarse links $\{0,1,5\}$, $\{1,2,4,5\}$, $\{2,3,4\}$, fine links
$\{5\}$, $\{4,5\}$, $\{2,3,4\}$, $\{1,2\}$, $\{0,1\}$, $r=0$, $s=1$
(\lean{not_theorem4LiteralND}). Verified for $m\ge4$ in every clause (\lean{theorem4_literal_ge4}).
\end{defect}

\begin{defect}{Bing--Jones 1959, Theorem~7, second sentence, p.~177}{False as printed for chains of
one or two links. Counted (one false statement). Ledger A.25 (third reading).}{1}{0}\label{it:bj7b}
``\emph{In fact, if $A$, $B$ are closed sets in $P\cdot d_1$, $P\cdot d_n$ respectively and
$H(x_r)=d_1$, $H(x_s)=d_n$, the chain $E$ may be selected so that $A\subset e_r$, $B\subset
e_s$.}'' With $n=1$, $m=3$, $H$ constant, $A=B=\{p\}$, $r=1$, $s=3$: $p\in e_1\cap e_3$, and $e_1$,
$e_3$ are not adjacent (\lean{bj_theorem7_ends_printed_false}). With $n=2$ the same, with $p\in
P\cdot d_1\cdot d_2$ and $H=(d_1,d_1,d_2)$ (a hand argument). For $n\ge3$ it holds when ``closed''
means closed in the space, which is Bing and Jones's own use on p.~178 (\lean{bj_theorem7_ends}).
Read as closed in the subspace $P\cdot d_1$, it fails for every $n$ (a hand argument:
$H=(d_1,d_2,d_1,d_2,d_3)$, $r=1$, $A=P\cdot d_1$).
\end{defect}

\subsection*{Tacit nonemptiness, not counted}
Jones 1948 Theorems~11--13 and 19 ($D$ or $U$ nonempty); van Mill 2004 Proposition~2.2 ($X\ne
\emptyset$; in the empty space $\emptyset$ is analytic and fat); Moore 1916's definition of
connected (the two parts must be nonempty); Moore 1925 Theorem~12 (a domain read as nonempty);
Cook 1967 Theorem~4 (the sequence $C_n$ does not exist when $C=X$, a case that needs no argument);
Whyburn (8.3) (``no convergent sequence of components'' must exclude the constant sequence).

%======================================================================
\section{Unargued steps and statements without proof}\label{sec:unargued}

Every step listed here is true and is supplied, or shown unnecessary, in the formalization,
except those marked ``not decided''. Sources are
listed in the order of \S\ref{sec:classification}.

\begin{defect}{Jones 1948, 1949, 1951 and 1955}{Unargued; stated without proof. Ledger A.10, A.11,
A.18.}{0}{0}\label{it:un-jones}
\begin{itemize}
\item Jones 1955, Theorem~1: the continuity step (``Since $G$ must be continuous at some one of its
  elements'': an upper semicontinuous decomposition of a compact metric space is continuous at a
  residual set of elements); homogeneity and aposyndesis of the decomposition space in (a); ``$N_o$
  is closed''. Not decided and moot: the attainment of the supremum $d(q)$ and ``$E_n$ is closed''.
\item Jones 1955, Theorem~2: ``if one element of $G$ were to separate $S$, then each of them
  would'' uses Borsuk's theorem; the ``uncountably many mutually exclusive domains'' (their choice
  and disjointness); the upper semicontinuity of $G'$ and ``$G$ is a compact subcontinuum of $G'$'';
  ``$M-g$ is connected''; the aposyndetic case ($L_x=\{x\}$) is not mentioned. Verified
  (\lean{jones55_theorem2'}). The three unnumbered claims of the closing remarks are verified.
\item Jones 1951, Theorem~1: ``$M'$ is homogeneous and aposyndetic'', ``$M'$ contains a
  nonseparating point'' and ``$M'$ is hereditarily unicoherent'' are asserted; the formal proof
  needs none of them.
\item Jones 1949: Case~1 (that every point of $M$ belongs to the boundary of an element of $G$ needs
  the Jordan curve theorem, \lean{jordan_curve_theorem}); four steps of Case~2, the last of which
  (``$M-(A+B)$ is \emph{not} connected'') needs $\theta$-curves and Janiszewski's theorem; and
  ``all except countably many \ldots\ must be of order 2'' needs $M$ uncountable.
\item Jones 1948: Theorems~0, 2, 3 and 5 are stated without proof (all proved); Theorem~13's choice
  of $V_2$ must satisfy $\overline{V_2}\subseteq W_1$.
\end{itemize}
\end{defect}

\begin{defect}{Whyburn 1939 \S14, Whyburn 1942, Moore 1922, Menger 1932}{Unargued; stated without
proof. Ledger A.19--A.21.}{0}{0}\label{it:un-why}
\begin{itemize}
\item Whyburn 1939, Theorem~14: ``$N$ can intersect only one of the arcs $xuy$, $xvy$''; ``the arc
  $x_1xuyy_1$ of $J$ is contained in $C$''; the limit continuum $H$ lies in $M$ and is
  nondegenerate; uniform continuity in the finite-cover step. Not decided, not needed: ``Hence
  there exists a simple closed curve $K$ in $V_\varepsilon(p)$ such that $K\cdot M=c+d$''. Corollary~1
  and (6.22) are stated without proof (supplied); footnote 16 (equivalence with Wilder's local
  0-avoidability) is not decided and not used.
\item Whyburn 1942: (8.4) ``there exist points $x$ and $y$ of $N$ \ldots\ such that
  $\delta[Q_2(x)\cdot Q_1(y)]<\varepsilon/4$'' needs facts from inside the proof of (2.2), not from
  the statement of (3.2); (8.41) uses the boundary bumping of (9.1); (9.2) ``and hence also of
  order 2 in $L$''. (3.1) is stated without proof (not used).
\item Moore 1922, Theorem~4: five steps of necessity and three of sufficiency, among them the sector
  theorem for continua crossing an annulus (proved by lifting to a strip) and the cyclic order of
  the continua $k_n$.
\item Menger 1932, p.~266: three steps of the \emph{Kennzeichnung der gewöhnlichen Kurven}; not
  used, since the formal proof follows Menger's second argument on p.~267.
\end{itemize}
\end{defect}

\begin{defect}{Hagopian 1975, 1976 and 1984}{Unargued. Ledger A.22--A.24.}{0}{0}\label{it:un-hag}
\begin{itemize}
\item Hagopian 1975: Lemma~1, three steps (the triod must be proper; complementary domains missing
  $M$ are pairwise disjoint; Baire and Moore~138); Theorem~1, paragraph 2 (``$F$, $f[F]$, and
  $g[F]$ are mutually disjoint''; ``$B\cup F\cup f[F]\cup g[F]$ is a triod''); paragraph 9, the
  second inclusion of ``$h_i[Y]\subset h_ik^{-1}[[r,t]]\subset E$''.
\item Hagopian 1976: Lemma~4, the second step (about two hundred lines of the formal proof,
  \lean{exists_proper_extension}); Lemmas~6, 7 and 9 (several sentences each); the plane
  Schoenflies theorem in Lemma~9 (\lean{plane_schoenflies}); ``we can assume \ldots\ that
  $k^{-1}[(0,1)]$ is accessible''; Properties 1 and 3; T7, T9, T11, T12, T13; ``It follows from
  Property 3 that $I$ is finite''; ``Note that $x_1=0$, $x_2=e$, and $x_{2n}=1$''; Property 4 for
  the top block; the order of the constants, which the text does not fix. Lemma~5's union of two
  extensions is not decided and not used. Every one of these is supplied except those listed in
  Item~\ref{it:h76}.
\item Hagopian 1984: Lemma~7 (three atriodicity steps, avoided); Lemma~8 (an arc); Lemma~12; Lemma~13
  (the hypotheses of Lemma~3, and an orientation argument for regularity); Lemmas~14--15
  (positivity of $\varepsilon=\frac12\rho(\bigcup\Delta(\mathcal P),M\setminus\bigcup\mathcal P')$);
  Lemma~16 (``Since $\mathcal P$ is free, there exist a 3-linked subchain \ldots'', replaced by a
  different lemma; Lemma~16 without normality is not decided); in Theorem~1, $X(Z)$, $J(Z)$ (without
  the printed positions of its end sets), claims (5)--(6), and ``$M/\mathfrak D$ \ldots\ contains an
  arc''.
\end{itemize}
\end{defect}

\begin{defect}{Jones 1975}{Unargued. Ledger A.28.}{0}{0}\label{it:un-j75}
(c) (the argument from (a) and (b) is not given). In the Theorem: ``let us assume that $T$ is
accessible from the complement of $M$''; ``There exists a positive integer $\eta$ \ldots\ This
follows from property (f)''; ``Clearly $t_f$ cannot lie in $\bar D_1\cup\bar D_2\cup\bar D_3$
without violating the selection of $\eta$''; ``it is easy to see that $p\notin\bar D_2$ and property
(f) is violated''. None is used by the formal proof (Item~\ref{it:j75}).
\end{defect}

\begin{defect}{Oversteegen--Tymchatyn 1982 and Cook 1967}{Unargued. Ledger A.29.}{0}{0}\label{it:un-ot}
\begin{itemize}
\item Theorem~3.6: ``Since $X$ is atriodic tree-like'' (not needed in the indecomposable case, so
  neither Rogers 1981 nor Jones 1969 is needed for Theorem~C); ``Hence there exists an index $i$
  and points \ldots\ such that $d(x,y)<\delta$''; ``since $Z$ is a proper subcontinuum''.
\item Corollary~3.5: ``For each $K\in\mathcal K$ we may suppose $K\cap\{\varepsilon\}\times
  R=\{x_K\}$'' (minimality does not give it in general; not needed); the uncountable family
  $\mathcal K$ of disjoint crossing continua.
\item Theorem~2.7: ``there exists a continuum $H$ in $C$ such that $(0,p)\in H$ and $H$ meets both
  $I\times\{0\}$ and $I\times\{1\}$'' and ``Clearly $K\cap H\ne\emptyset$''; the continuity of
  $\varphi$ can fail where the path meets $B_1$ (harmless where it is used). Remark~2.5 and the
  sketch of Lemma~2.1.
\item Cook 1967, Theorem~4: the sequence $C_n$ is not constructed, and ``The common part of
  $f(C_1)$, $f(C_2)$, \ldots\ is $f(C)$'' is unstated; the formal proof needs neither
  (\lean{cook67_theorem4}).
\end{itemize}
\end{defect}

\begin{defect}{Hoehn--Oversteegen 2016}{Unargued; stated without proof. Ledger A.26.}{0}{0}\label{it:un-ho}
\begin{itemize}
\item p.~211, proof of Theorem~1: ``It follows that $d_H(X,T)<\delta$''; from $d_{\sup}(f,
  \mathrm{id}_X)<\delta$ only $X\subseteq T_\delta$ and $f(X)\subseteq X_\delta$ follow, while
  Theorem~4 needs $T\subseteq X_\delta$ too. Repair: replace $T$ by $f(X)$.
\item Theorem~8, (1)$\Rightarrow$(2): the fans must be pairwise disjoint, which is not said.
\item Lemma~6 (p.~185): ``The proofs of these properties are left to the reader''; seven properties,
  verified (\lean{ho16_lemma6}).
\item pp.~186--187: the fold determined by $G_1$, $G_2$, $G_3$ (``It is straightforward to see'') and
  Proposition~7 (three steps, supplied).
\item pp.~189--191: the observations after Definition~9; Proposition~11 (three short steps).
\item pp.~194--198, Theorem~15: Claim~15.1 (``We leave it to the reader''); the passage to an
  irreducible subset; the branch-point surgery; ``clearly $S$ irreducibly separates''; ``(S3) and
  (S5) are immediate''. The formal proof avoids all of them (\lean{ho16_theorem15}).
\item Lemma~18 (that $\pi(S')$ has finitely many components); Proposition~17 (``It is now easy to
  verify''); pp.~203--210 (the ``straightforward observations'', (S3) and (S5$'$) in
  Proposition~19, the component $K$ of p.~209).
\end{itemize}
\end{defect}

\begin{defect}{Bing--Jones 1959}{Unargued. Ledger A.25.}{0}{0}\label{it:un-bj}
\begin{itemize}
\item p.~173: ``It follows from Theorem 10 that any two circles of pseudo-arcs are topologically
  equivalent and that each is homogeneous''; p.~189: ``It follows from Theorem 10 that each circle
  of pseudo-arcs is homogeneous and that any two of them are homeomorphic.'' The passage from arcs
  to circles is not written; it needs circle-likeness to make the two halves snake-like
  (\lean{isSnakeLike_of_subset_circleLike}), and then Theorem~10 and a gluing.
\item Theorem~2 ($G(b,c)$ meets $d_t$); Theorem~3 ($D$ an $\varepsilon$-chain); Theorem~7 (above);
  Theorem~8 ($U_A$, $U_B$, asserted for an upper semicontinuous collection; the cover
  $U_1,\ldots,U_t$, ``It follows from Theorem 4''; the general case, ``alter $Y$ as was done in the
  proof of Theorem 7''); the corollary to Theorem~8 (three consecutive links); Theorem~9 (``That $h$
  is a homeomorphism of $P$ onto $Q$ follows from an argument similar to'').
\item Theorem~10: Step~1(b); Step~2(b), the chain map $R_{i/n}$; Step~2(c), Theorem~6 applied
  without naming the chain in the role of $E$, the use of the false Theorem~5
  (Item~\ref{it:bj5}), and the transition set the repaired Theorem~5 needs; ``Of course,
  $d(p,i+1)\subset d(p,i)$'' (p.~188, from Condition B); Step~2(i), Theorem~9's hypothesis; Step~3 (closures of the $f_j$); ``Steps
  4, 5, $\cdots$. These steps are essentially repetitions of Step 3.'' Verified
  (\lean{bjThm10}, through \lean{bj_step}).
\end{itemize}
\end{defect}

\begin{defect}{Bing 1948 and 1951; Fugate 1966 and 1969; Lehner 1961; van Mill 2004}{Unargued; stated
without proof; not established for the generality stated. Ledger A.1--A.9.}{0}{0}\label{it:un-bing}
\begin{itemize}
\item Bing 1948, Theorem~14: ``That there exist such sequences follows as a consequence of Theorem
  9'' is not a proof; the descriptions are supplied by Lehner's Lemma~3 for $n=2$
  (\lean{lehner_lemma3_two}). Two assertions of the sketch are not established for an
  \emph{arbitrary} covering description; both hold for the description Lehner's construction
  produces. Nine further steps are supplied.
\item Bing 1951 (Duke): Theorem~7 is stated for tree-like continua, but the proof of Theorem~6 that
  it rests on assumes that no subcontinuum separates the plane; ``tree-like $\Rightarrow$ no
  separating subcontinuum'' is not stated (verified, \lean{treeLike_nonseparating}). Theorem~8:
  ``$K_2$ contains $K_4$ interiorly'' and the order of the two inclusions (\lean{theorem8_1951'}).
  Theorem~11, sufficiency: three sentences, supplied from Fugate 1969 (\lean{theorem11_1951}).
\item Fugate 1966, Corollary~1: stated without proof; the induction it suggests needs every partial
  union to be unicoherent, which is what is missing; verified (\lean{fugate_cor1},
  \lean{fugateU_holds}). No proof or discussion of the corollary for more than two pieces was found
  in the reviews, in later work of Fugate, in Nadler 1992, Macías 2005, the Houston Problem Book,
  or the citing papers searched.
\item Fugate 1969: besides Lemma~3 (\S\ref{sec:local}), nine smaller steps (hereditary unicoherence
  not stated; the displayed chain is only a weak chain).
\item van Mill 2004, Proposition~3.5: Proposition~2.2 is applied ``in $E$'', which needs ``nowhere
  meager in $X$'' to give ``nowhere meager in $E$''; avoided by a relative form of 2.2.
\end{itemize}
\end{defect}

\begin{defect}{Moore 1916, 1919 and 1925}{Unargued; stated without proof. Ledger
A.12--A.16.}{0}{0}\label{it:un-moore}
\begin{itemize}
\item Moore 1916: Theorems~1, 3, 7, 11, 13, 16--18, 20, 22, 30, 38, 40, 42 and 43 are stated without
  proof (all proved; Theorem~1 needs no axiom); \S5 rests on Lennes 1911 and supplies no proofs;
  Theorem~41's last step (``It may be proved that $J$ is a closed curve satisfying the
  requirements'') hides a limit construction; Theorems~37 and 39 (about seven steps). The opening
  of \S6 (``It is clear that Theorems 1--19 are consequences of Axioms 2--5 and Theorems 5 and 15'')
  is partly borne out and not formalized.
\item Moore 1919: Theorem~B (five steps called ``clear'' or ``easy''); Theorems~D and E (``This
  process may be continued'', and the extension of rulings, never proved); Theorem~E's limit arcs;
  the nine-cell ruling of Fig.~3 (not needed); Theorem~G (``It is clear that through every point
  there is just one curve of the set $\alpha$ \ldots''); the final paragraph (``It is clear that the
  correspondence thus established \ldots\ is continuous''). Verified by a route through a Schoenflies
  theorem for Moore's space (\lean{m19_theoremG}, \lean{m19_plane}). The claim of the introduction
  for $\Sigma_2$ rests on Moore 1916 p.~163 (Item~\ref{it:m25-21}).
\item Moore 1925: Theorem~2 (Brouwer's theorem is needed twice, and three ``clearly'' steps);
  Theorem~11 (``therefore neither $\bar a$ nor $\bar b$ separates $S$''; ``It is easy to see that
  one of these domains \ldots\ is bounded''; ``Clearly $I$ and $E$ are domains \ldots\ and
  $I+E=G-J$'');
  Theorem~3 (four steps); Theorems~4, 5, 6, 9, 10 and 15--19 and the p.~420 claim, stated without
  proof (supplied); Theorems~8, 13, 14 and 20 (two to three steps each).
\end{itemize}
\end{defect}

%======================================================================
\section{Misprints and harmless mismatches}\label{sec:misprints}

Each is read as indicated; none affects a result. ``Printed $X$, read $Y$.''

\begin{itemize}
\item \textbf{Hoehn--Oversteegen 2016.} p.~211, Theorem~1 restated without ``non-degenerate''
  (Item~\ref{it:ho1p}); p.~189, Theorem~8, (3)$\Rightarrow$(1): $U$ and $V$ interchanged in the first
  sentence; p.~195: $x$ for $(x,y)$; p.~198: ``Since $a_i\in C$'' for $a_p$, and $\{z\}_*$ for the
  star; p.~202: $\cap$ for $\cup$ in ``$\mathcal E(S'')\subset(F_1\cap F_2)_\star$'', and ``$S''$ is a
  straight'' (a word missing); pp.~205--207: $\varphi^{-1}$ for $\varphi_\star^{-1}$, three times,
  and on p.~205 ``$P'_1\cap(\partial F_3)_\star=\varphi_\star^{-1}(\gamma_{i_0})$'' for its
  intersection with $(\partial F_3)_\star$; p.~209: the order of composition; p.~192:
  ``comple ment''; p.~210 says more than p.~201; Theorem~20 (p.~210): an unneeded assumption on the
  $\varepsilon_i$; pp.~202--203: two omitted citations of (S4$'$).
\item \textbf{Bing--Jones 1959.} p.~172: \emph{of} for \emph{if}; p.~173: \emph{of} for
  \emph{are}; Theorem~3: $F(v,m)$ for $F(w,m)$; Theorem~7: ``components'' for ``composants'';
  Theorem~8: $d_2$ for $d_n$; Theorem~9: $G$ for $F$, and $e$ for $d$ on p.~180; p.~181: $E_i$ for
  $E_1$, $H_2$ for $H_1$, $d^1_k$ for $e^1_k$; p.~182: $K_3$ for $K_2$; p.~185:
  $R_{(2i+1)/n}$ for $R_{(2i+1)/2n}$ and $P(0,1/n-\varepsilon)$ for $P(0,i/n-\varepsilon)$; p.~186:
  $E_k$ for $E_1$; p.~187: ``intesects'', $E_1$ for $D_1$, $f_{i+1}$ for $f_{t+1}$; p.~178, the
  corollary: $4i+1,\ldots,4i+4$ for $4i-3,\ldots,4i$; p.~188: $H_i(d_k^{i+1})$; p.~189:
  $\rho(p,q)<4/2^i$ where the meshes give $8/2^i$; p.~180: the last display of Theorem~9 ends
  ``$\subset\bar e_i$'' where the claim is $h(P\cdot d_i)\subset e_i$; p.~187: condition (1) asks
  for the whole link $f_j$ where the construction controls only $f_j\cdot Q(i/k,(i+1)/k)$; p.~191, the figure: the level written second,
  and $a_1$, $c_1$ for $a_i$, $c_i$; pp.~190--191: the letter $G$ used twice.
\item \textbf{Hagopian.} 1975: ``transfromation'' (p.~37); Lemma~3 printed for type A, citing a
  source not in hand (only type A$'$ is used). 1976: none beyond the citations of
  Item~\ref{it:hag-cite}. 1984: ``opens sets'' (p.~333), ``$\varepsilon$-circluar'' (p.~334),
  ``$\varepsilon$-homomorphisms'' (p.~338), unbalanced brackets in Lemma~13(5), Lemma~13(5)
  printed without $h$ (p.~340), $K(Z)$ ``the component of $\bigcup\mathcal P'$ that contains $Z$''
  (read: whose closure has $Z$ as a layer); Lemma~2 (cited to Mackowiak--Tymchatyn 13.8) is
  applied to unions $F\cup J$, $N\cup F$, which is harmless in the plane.
\item \textbf{Oversteegen--Tymchatyn 1982.} ``$\sigma_0(x)$'' (p.~129); ``$f(X)\supset g(X)$'' for
  $f(C)\supset g(C)$ and ``homomorphism'' (p.~140); ``[7]'' for Effros, which is [6]
  (pp.~139--140); $R^2$ for $I^2$ (p.~136); ``$d(h_\lambda(Z),h_\gamma(z))<1/8$'' for
  $h_\lambda(z)$ (p.~138); ``an indecomposable homogeneous plane'' (``continuum'' missing).
\item \textbf{Jones.} 1948, Theorem~16: $P$ for the weak cut point $p$. 1951: ``hereditorily''
  (p.~858). 1955, Theorem~2's proof: $g$ for $G$.
\item \textbf{Moore.} 1916: p.~139, ``between $P_1$ and $B$'' must allow $B$; Theorem~33, an
  unbalanced parenthesis; Theorem~34, ``Jordan \emph{arc}'' for ``closed curve''. 1919: footnote
  $\dagger$ (p.~170), $A$ and $B$ for $X$ and $Y$. 1922: Theorem~1, $K^*$ for $\bar K$; Theorem~3,
  ``diameter less than $\varepsilon$'' where only $\le2\varepsilon$ holds (harmless with
  $\varepsilon/3$); Theorem~4, $t_m$ for $\bar t_m$, $1\le m\le n$ for $4^n$, and
  $\varepsilon=(b-c)/4$ for $(c-b)/4$; ``fie'', ``curce'', ``a sumed''. 1925: Theorem~1 (two slips);
  Theorem~14, $pxq$ with $x$ never introduced; Theorem~18, $x$ and $y$ not introduced.
\item \textbf{Whyburn 1939.} ``$ou-o$'' for $ou-u$; ``$<3\varepsilon$'' after ``$<4\varepsilon$''
  (the argument gives $2\varepsilon$); ``there exists points''.
\item \textbf{Fugate 1966.} The introduction says ``a-triodic and hereditarily unicoherent'',
  Theorem~1 ``a-triodic and unicoherent'' (equivalent for two pieces).
\item \textbf{Domokos--Prajs 2022.} ``homogenous'' (p.~13).
\end{itemize}

%======================================================================
\appendix
\section{Readings refuted that are not defects of the sources}\label{app:readings}

Each of the following is a reading of a printed statement that the formalization refuted, and that
the source does not make. They are recorded so that the refutations are not mistaken for defects.

\begin{defect}{Bing 1951 (Duke), Theorem~8, p.~658, with $g_p$ read with one layer}{A misreading,
not a defect: Theorem~8 is true as printed. Not counted. Ledger A.2.1, B.1.}{0}{0}\label{it:b51-8}
Bing defines (p.~658): ``For each point $p$ let $g_p$ denote the intersection of all continua
which contain interiorly a continuum that contains $p$ interiorly.'' With the inner continuum
dropped, $g_p$ is the intersection of the continua containing $p$ interiorly, and Theorem~8 fails:
in the plane continuum $S_0$ made of two sine curves over a bar, $q\in g_p$ while $g_p\ne g_q$
(\lean{Counter92.theorem8_1951_false}). With Bing's definition every clause is verified
(\lean{theorem8_1951'}), clause (a), ``$G$ is an arc with respect to its elements'', in the form
Bing's proof establishes: of any three elements, one separates the other two; on $S_0$ the elements are the bar and the points of the curves
(\lean{Counter92.bingElt_bar}, \lean{Counter92.bingElt_curve}). The inner layer is what makes the
elements saturate.
\end{defect}

\begin{defect}{Other refuted readings}{Not defects. Ledger B.1.}{0}{0}\label{it:readings}
\begin{itemize}
\item Bing 1948, Theorem~3, end clause read ``respectively'': false at $k=3$
  (\lean{Recon121.not_theorem3Literal}); Bing's own ``or'' is proved.
\item ``1951-crooked implies 1948-crooked'': the two papers' notions of crookedness differ
  (\lean{isCrookedIn_ne_isCrookedIn'}; four strips and six rectangles).
\item Bing 1951, Theorem~11, the separation read with $K'\supseteq K-g'-N$: refuted.
\item Nadler, \emph{Continuum Theory}, 11.28 read without $U-H\ne\emptyset$: false for $H=\bar U$
  (\lean{isConnected_diff_of_superset_false}); Nadler's proof assumes it.
\item van Mill 2004, Lemma~3.4 and Proposition~3.5 read without their closure bars: the printed
  statements, with the bars, are right.
\end{itemize}
\end{defect}

\section{Lean names}\label{app:lean}
Every name cited in the text, with its full name and the axioms its audit lists in
\lean{axioms.txt}. ``Std'' is \lean{propext}, \lean{Classical.choice}, \lean{Quot.sound};
a dash means no axioms. ``Used'' marks the declarations that the two final theorems use
(the dependency census, \lean{doc/depcensus/used.txt}). The last column gives the items
that cite the name.
\begin{small}
\begin{longtable}{@{}>{\raggedright\arraybackslash}p{0.60\textwidth}>{\raggedright\arraybackslash}p{0.09\textwidth}>{\raggedright\arraybackslash}p{0.08\textwidth}>{\raggedright\arraybackslash}p{0.15\textwidth}@{}}
\toprule Name & Axioms & Used & Items\\ \midrule \endhead
\lean{HPC.BJ59.bj_step} & Std & yes & 46\\
\lean{HPC.BJ59.bj_theorem3} & Std &  & \S7\\
\lean{HPC.BJ59.bj_theorem5} & Std &  & 8\\
\lean{HPC.BJ59.bj_theorem5_printed_false} & Std &  & 8\\
\lean{HPC.BJ59.bj_theorem5_printed_false_general} & Std &  & 8\\
\lean{HPC.BJ59.bj_theorem7} & Std &  & 19\\
\lean{HPC.BJ59.bj_theorem7_ends} & Std &  & 39\\
\lean{HPC.BJ59.bj_theorem7_ends_printed_false} & Std &  & 39\\
\lean{HPC.BJ59.bj_theorem7_proof_n3} & Std &  & 19\\
\lean{HPC.BJ59.bj_theorem8} & Std & yes & \S7\\
\lean{HPC.BJ59.bj_theorem9} & Std &  & \S7\\
\lean{HPC.BJ59.bjPlaneExample} & Std & yes & 2\\
\lean{HPC.BJ59.bjThm10} & Std & yes & 1, 46\\
\lean{HPC.BJ59.circ_step} & Std & yes & 2\\
\lean{HPC.BJ59.CircDescr.isCircleLike_M} & Std & yes & 2\\
\lean{HPC.BJ59.CircDescr.isCircleOfPA} & Std & yes & 2\\
\lean{HPC.BJ59.dp22_theorem61'} & Std &  & 1\\
\lean{HPC.BJ59.dp22_theorem61_iff'} & Std &  & 1, 14\\
\lean{HPC.BJ59.exists_circDescr_plane} & Std & yes & 2\\
\lean{HPC.BJ59.ho16_classification_plane₄} & Std & yes & \S2, 1\\
\lean{HPC.BJ59.ho16_classification₄} & Std & yes & \S2, 1\\
\lean{HPC.BJ59.hoCircleLike'} & Std &  & 1\\
\lean{HPC.BJ59.isSnakeLike_of_atomic'} & Std & yes & 1\\
\lean{HPC.BJ59.sweep_lemma} & Std & yes & 2\\
\lean{HPC.BJ59.sweep_lemma3} & Std &  & 2\\
\lean{HPC.BJ59.theorem5_repaired} & Std & yes & 8\\
\lean{HPC.Fold.not_theorem4Literal} & Std &  & 38\\
\lean{HPC.Fold.not_theorem4LiteralND} & Std &  & 38\\
\lean{HPC.Fold.not_theorem8_literal} & Std &  & 9\\
\lean{HPC.Fold.Recon121.not_theorem3Literal} & Std &  & 50\\
\lean{HPC.Fold.Recon121.seam_every_H_fails} & Std &  & 6\\
\lean{HPC.Fold.Recon121.seamD_special} & Std &  & 6\\
\lean{HPC.Fold.Recon121.theorem3LiteralOr_holds} & Std &  & 6\\
\lean{HPC.Fold.Recon122.bing_main} & Std &  & 6\\
\lean{HPC.Fold.Recon122.special_five} & Std &  & 6\\
\lean{HPC.Fold.theorem3_bing} & Std &  & 6\\
\lean{HPC.Fold.theorem3_literal} & Std &  & 6\\
\lean{HPC.Fold.theorem4_literal_ge4} & Std &  & 38\\
\lean{HPC.Fold.theorem4_two_degenerate_false} & Std &  & 38\\
\lean{HPC.Fold.theorem7_minimal} & Std &  & 9\\
\lean{HPC.Fold.theorem8} & Std &  & 9\\
\lean{HPC.Fold.theorem9} & Std &  & 9\\
\lean{HPC.Hagopian.bing60_plane_uncond} & Std &  & 4\\
\lean{HPC.Hagopian.blockLayersMeetSides3} & Std &  & 3\\
\lean{HPC.Hagopian.blockLayersMeetSides_false} & Std &  & 3\\
\lean{HPC.Hagopian.exists_proper_extension} & Std & yes & 42\\
\lean{HPC.Hagopian.exists_separating_subcontinuum_of_homeomorph} & Std & yes & 17\\
\lean{HPC.Hagopian.hagopian76_far_layer} & Std &  & \S7\\
\lean{HPC.Hagopian.hagopian76_lemma4_layers'} & Std &  & \S7\\
\lean{HPC.Hagopian.hagopian76_theorem1} & Std & yes & 3\\
\lean{HPC.Hagopian.hagopian76_theorem1_of_cases4} & Std &  & 3\\
\lean{HPC.Hagopian.hagopian84_cover} & Std & yes & \S7\\
\lean{HPC.Hagopian.hagopian84_decomp} & Std & yes & 3\\
\lean{HPC.Hagopian.hagopian84_lemma15} & Std & yes & \S7\\
\lean{HPC.Hagopian.hagopian84_lemma7} & Std &  & 3\\
\lean{HPC.Hagopian.innerComponentsCross_hagopian} & Std &  & 3\\
\lean{HPC.Hagopian.jones75_contra} & Std & yes & 4\\
\lean{HPC.Hagopian.jones75_noArc} & Std & yes & 4\\
\lean{HPC.Hagopian.JonesSetting.jones75_d} & Std &  & 4\\
\lean{HPC.Hagopian.JonesSetting.nested} & Std & yes & 4\\
\lean{HPC.Hagopian.JonesSetting.orientation} & Std & yes & 4\\
\lean{HPC.Hagopian.JonesSetting.union_arc} & Std & yes & 4\\
\lean{HPC.Hagopian.last_step_ending} & Std &  & 3\\
\lean{HPC.Hagopian.no_proper_subcontinuum_separates} & Std & yes & 17\\
\lean{HPC.Hagopian.plane_schoenflies} & Std &  & 42\\
\lean{HPC.Hagopian.sUnion_cover84_eq_univ} & Std &  & \S7\\
\lean{HPC.Hagopian.trap_jordan} & Std & yes & 4\\
\lean{HPC.HO16.ho16_indec_homogeneous_iff} & Std & yes & \S2\\
\lean{HPC.HO16.ho16_lemma6} & Std & yes & 45\\
\lean{HPC.HO16.ho16_p193_remark_false} & Std &  & \S8\\
\lean{HPC.HO16.ho16_p208_false} & Std &  & 20\\
\lean{HPC.HO16.ho16_prop10} & Std & yes & \S7\\
\lean{HPC.HO16.ho16_prop19} & Std & yes & 20\\
\lean{HPC.HO16.ho16_prop19_full} & Std & yes & 20\\
\lean{HPC.HO16.ho16_theorem15} & Std & yes & \S7, 45\\
\lean{HPC.HO16.ho16_theorem1_plane} & Std & yes & \S2\\
\lean{HPC.HO16.ho16_theorem1_printed_false} & Std &  & 37\\
\lean{HPC.HO16.ho16_theorem1prime_printed_false} & Std &  & 37\\
\lean{HPC.HO16.ho16_theorem20} & Std & yes & \S2\\
\lean{HPC.isCrookedIn_ne_isCrookedIn'} & Std &  & 50\\
\lean{HPC.Jones.aposAt_of_isContinuum_interior} & Std & yes & 5\\
\lean{HPC.Jones.jones49_order_two} & Std & yes & 18\\
\lean{HPC.Jones.jones49_theorem1} & Std & yes & 36\\
\lean{HPC.Jones.jones49_theorem2} & Std &  & 36\\
\lean{HPC.Jones.jones55_a} & Std & yes & 5\\
\lean{HPC.Jones.jones55_b} & Std & yes & 5\\
\lean{HPC.Jones.jones55_c} & Std & yes & 1, 5, 14\\
\lean{HPC.Jones.jones55_continuous} & Std & yes & 5\\
\lean{HPC.Jones.jones55_elements} & Std & yes & 5\\
\lean{HPC.Jones.jones55_theorem2'} & Std & yes & 40\\
\lean{HPC.Jones.jones55_usc} & Std & yes & 5\\
\lean{HPC.Jones.jordan_curve_theorem} & Std & yes & 40\\
\lean{HPC.Jones.moore22_theorem4} & Std &  & 27\\
\lean{HPC.Jones.not_jones49_theorem1_point} & Std &  & 36\\
\lean{HPC.Jones.not_jones49_theorem2_point} & Std &  & 36\\
\lean{HPC.Jones.not_theorem10_printed_unitInterval} & Std &  & 33\\
\lean{HPC.Jones.not_theorem6_printed} & Std &  & 33\\
\lean{HPC.Jones.theorem10_printed_iff} & Std &  & 33\\
\lean{HPC.Jones.whyburn39_theorem14} & Std & yes & 11\\
\lean{HPC.Jones.whyburn_9_1} & Std & yes & 28\\
\lean{HPC.lehner_lemma3_two} & Std &  & 47\\
\lean{HPC.Moore.fa_theorem10} & Std &  & 21\\
\lean{HPC.Moore.fa_theorem14} & Std &  & 34\\
\lean{HPC.Moore.fa_theorem15} & Std &  & 22\\
\lean{HPC.Moore.fa_theorem19} & Std & yes & 34\\
\lean{HPC.Moore.fa_theorem23} & Std &  & 34\\
\lean{HPC.Moore.fa_theorem28} & Std & yes & 23\\
\lean{HPC.Moore.fa_theorem33} & Std & yes & 24\\
\lean{HPC.Moore.fa_theorem34} & Std &  & 24\\
\lean{HPC.Moore.fa_theorem35} & Std & yes & 34\\
\lean{HPC.Moore.fa_theorem36} & Std & yes & 25\\
\lean{HPC.Moore.fa_theorem36_interior_subset} & Std & yes & 25\\
\lean{HPC.Moore.fa_theorem4} & Std & yes & 34\\
\lean{HPC.Moore.fa_theorem43} & Std &  & 34\\
\lean{HPC.Moore.frontier_eq_of_scc} & Std & yes & 26\\
\lean{HPC.Moore.m19_plane} & Std & yes & 10, 48\\
\lean{HPC.Moore.m19_theoremC_corrected} & Std &  & 10\\
\lean{HPC.Moore.m19_theoremC_crosscuts} & Std &  & 10\\
\lean{HPC.Moore.m19_theoremC_within} & Std &  & 10\\
\lean{HPC.Moore.m19_theoremD} & Std & yes & 10\\
\lean{HPC.Moore.m19_theoremG} & Std &  & 48\\
\lean{HPC.Moore.m25_theorem10} & Std & yes & 30\\
\lean{HPC.Moore.m25_theorem11} & Std & yes & 26, 30\\
\lean{HPC.Moore.m25_theorem2} & Std & yes & 29\\
\lean{HPC.Moore.m25_theorem21} & Std & yes & 12\\
\lean{HPC.Moore.m25_theorem22} & Std & yes & 12\\
\lean{HPC.Moore.m25_theorem7} & Std &  & 31\\
\lean{HPC.Moore.m25_unbounded_unique} & Std & yes & 29\\
\lean{HPC.Moore.mullikin_theorem4_disjoint} & Std &  & \S8\\
\lean{HPC.Moore.mullikin_theorem5} & Std & yes & 32\\
\lean{HPC.Moore.not_fa_sec5_limitPt_printed} & Std &  & 34\\
\lean{HPC.Moore.not_fa_theorem14_printed} & Std &  & 34\\
\lean{HPC.Moore.not_fa_theorem43_printed} & Std &  & 34\\
\lean{HPC.Moore.not_m19_theoremC_printed} & Std &  & 10\\
\lean{HPC.Moore.not_m25_theorem10_printed} & Std &  & 30\\
\lean{HPC.Moore.not_m25_theorem11_printed} & Std &  & 30\\
\lean{HPC.Moore.not_m25_theorem11_singleton} & Std &  & 35\\
\lean{HPC.Moore.not_m25_theorem2_printed} & Std &  & 29\\
\lean{HPC.Moore.not_m25_theorem7_printed} & Std &  & 31\\
\lean{HPC.Moore.not_moore1920_theorem4_singleton} & Std &  & 35\\
\lean{HPC.Moore.not_mullikin_theorem4_disjoint_mathlib} & Std &  & \S8\\
\lean{HPC.Moore.not_mullikin_theorem5_printed} & Std &  & 32\\
\lean{HPC.Plane.Counter92.bingElt_bar} & Std &  & 49\\
\lean{HPC.Plane.Counter92.bingElt_curve} & Std &  & 49\\
\lean{HPC.Plane.Counter92.theorem8_1951_false} & Std &  & 49\\
\lean{HPC.Plane.fugate_cor1} & Std &  & 47\\
\lean{HPC.Plane.fugateU_holds} & Std &  & 47\\
\lean{HPC.Plane.isConnected_diff_of_superset_false} & Std &  & 50\\
\lean{HPC.Plane.isHereditarilyUnicoherent_union} & Std &  & 16\\
\lean{HPC.Plane.isSnakeLike_of_subset_circleLike} & Std & yes & 46\\
\lean{HPC.Plane.miller_1_5_false} & Std &  & 16\\
\lean{HPC.Plane.theorem11_1951} & Std &  & 47\\
\lean{HPC.Plane.theorem8_1951'} & Std &  & 47, 49\\
\lean{HPC.Plane.treeLike_nonseparating} & Std &  & 47\\
\lean{HPC.Span.cook67_theorem4} & Std & yes & 44\\
\lean{HPC.Span.ot82_cor35_of_thm32} & Std & yes & 7\\
\lean{HPC.Span.ot82_theorem32} & Std & yes & 7\\
\lean{HPC.Span.ot82_theorem36_of_cor35} & Std & yes & 7\\
\lean{HPC.Span.theoremC} & Std & yes & 7\\
\bottomrule
\end{longtable}
\end{small}

%======================================================================
\begin{thebibliography}{HMP20}\small
\bibitem[B48]{B48} R. H. Bing, A homogeneous indecomposable plane continuum, \emph{Duke Math. J.}
  15 (1948), 729--742.
\bibitem[B51a]{B51a} R. H. Bing, Snake-like continua, \emph{Duke Math. J.} 18 (1951), 653--663.
\bibitem[B51b]{B51b} R. H. Bing, Concerning hereditarily indecomposable continua, \emph{Pacific J.
  Math.} 1 (1951), 43--51.
\bibitem[B60]{B60} R. H. Bing, A simple closed curve is the only homogeneous bounded plane continuum
  that contains an arc, \emph{Canad. J. Math.} 12 (1960), 209--230.
\bibitem[BJ59]{BJ59} R. H. Bing and F. B. Jones, Another homogeneous plane continuum, \emph{Trans.
  Amer. Math. Soc.} 90 (1959), 171--192.
\bibitem[C67]{C67} H. Cook, Continua which admit only the identity mapping onto non-degenerate
  subcontinua, \emph{Fund. Math.} 60 (1967), 241--249.
\bibitem[DP22]{DP22} A. Domokos and J. R. Prajs, The Planck boundary within the hyperspace of the
  circle of pseudo-arcs, arXiv:2203.15906 (2022).
\bibitem[F66]{F66} J. B. Fugate, Decomposable chainable continua, \emph{Trans. Amer. Math. Soc.}
  123 (1966), 460--468.
\bibitem[F69]{F69} J. B. Fugate, A characterization of chainable continua, \emph{Canad. J. Math.} 21
  (1969), 383--393.
\bibitem[H75]{H75} C. L. Hagopian, Homogeneous plane continua, \emph{Houston J. Math.} 1 (1975),
  35--41.
\bibitem[H76]{H76} C. L. Hagopian, Indecomposable homogeneous plane continua are hereditarily
  indecomposable, \emph{Trans. Amer. Math. Soc.} 224 (1976), 339--350.
\bibitem[H84]{H84} C. L. Hagopian, Atriodic homogeneous continua, \emph{Pacific J. Math.} 113
  (1984), 333--347.
\bibitem[HMP20]{HMP20} C. L. Hagopian, M. M. Marsh and J. R. Prajs, Folders of continua,
  \emph{Topology Proc.} 55 (2020), 13--33.
\bibitem[HO16]{HO16} L. C. Hoehn and L. G. Oversteegen, A complete classification of homogeneous
  plane continua, \emph{Acta Math.} 216 (2016), 177--216.
\bibitem[I68]{I68} W. T. Ingram, Decomposable circle-like continua, \emph{Fund. Math.} 63 (1968),
  193--198.
\bibitem[J48]{J48} F. B. Jones, Concerning non-aposyndetic continua, \emph{Amer. J. Math.} 70
  (1948), 403--413.
\bibitem[J49]{J49} F. B. Jones, Note on homogeneous plane continua, \emph{Bull. Amer. Math. Soc.} 55
  (1949), 113--114.
\bibitem[J51]{J51} F. B. Jones, Certain homogeneous unicoherent indecomposable continua,
  \emph{Proc. Amer. Math. Soc.} 2 (1951), 855--859.
\bibitem[J55]{J55} F. B. Jones, On a certain type of homogeneous plane continuum, \emph{Proc. Amer.
  Math. Soc.} 6 (1955), 735--740.
\bibitem[J75]{J75} F. B. Jones, Use of a new technique in homogeneous continua, \emph{Houston J.
  Math.} 1 (1975), 57--61.
\bibitem[L61]{L61} G. R. Lehner, Extending homeomorphisms on the pseudo-arc, \emph{Trans. Amer.
  Math. Soc.} 98 (1961), 369--394.
\bibitem[Le79]{Le79} A. Lelek, The span of mappings and spaces (research announcement),
  \emph{Topology Proc.} 4 (1979), 631--633.
\bibitem[Lw83]{Lw83} W. Lewis, Homogeneous continua and continuous decompositions, \emph{Topology
  Proc.} 8 (1983), 71--84.
\bibitem[Lw85]{Lw85} W. Lewis, Continuous curves of pseudo-arcs, \emph{Houston J. Math.} 11 (1985),
  91--99.
\bibitem[Lw96]{Lw96} W. Lewis, Continuous collections of hereditarily indecomposable continua,
  \emph{Topology Appl.} 74 (1996), 169--176.
\bibitem[M16]{M16} R. L. Moore, On the foundations of plane analysis situs, \emph{Trans. Amer. Math.
  Soc.} 17 (1916), 131--164.
\bibitem[M19]{M19} R. L. Moore, Concerning a set of postulates for plane analysis situs,
  \emph{Trans. Amer. Math. Soc.} 20 (1919), 169--178.
\bibitem[M20]{M20} R. L. Moore, Concerning simple continuous curves, \emph{Trans. Amer. Math. Soc.}
  21 (1920), 333--347.
\bibitem[M22]{M22} R. L. Moore, Concerning connectedness im kleinen and a related property,
  \emph{Fund. Math.} 3 (1922), 232--237.
\bibitem[M25]{M25} R. L. Moore, Concerning upper semi-continuous collections of continua,
  \emph{Trans. Amer. Math. Soc.} 27 (1925), 416--428.
\bibitem[Mu22]{Mu22} A. M. Mullikin, Certain theorems relating to plane connected point sets,
  \emph{Trans. Amer. Math. Soc.} 24 (1922), 144--162.
\bibitem[OT82]{OT82} L. G. Oversteegen and E. D. Tymchatyn, Plane strips and the span of continua
  (I), \emph{Houston J. Math.} 8 (1982), 129--142.
\bibitem[P00]{P00} J. R. Prajs, A continuous circle of pseudo-arcs filling up the annulus,
  \emph{Trans. Amer. Math. Soc.} 352 (2000), 1743--1757.
\bibitem[vM04]{vM04} J. van Mill, A note on the Effros theorem, \emph{Amer. Math. Monthly} 111
  (2004), 801--806.
\bibitem[W39]{W39} G. T. Whyburn, Semi-locally-connected sets, \emph{Amer. J. Math.} 61 (1939),
  733--749.
\bibitem[W42]{W42} G. T. Whyburn, \emph{Analytic Topology}, Amer. Math. Soc. Colloq. Publ. 28,
  1942.
\end{thebibliography}

\section*{Acknowledgments}
The formalization was carried out by the Aristotle system (Harmonic) in Lean~4 with Mathlib, from
commissions written by Claude (Anthropic), who also compiled the defects ledger and this document
from it.

\end{document}
